\chapter{多模态生成}
\label{chap:vision-multimodal-generation}

“让玩具狗绕花瓶走一圈，碰到花瓶时发出轻响”看似只有一句要求，交付时却要同时回答几个问题。玩具狗在转身后是否仍是原来的样子，花瓶有没有真的受到碰撞，声音是否发生在接触时刻？如果还要把玩具狗做成可以转动观察的三维资产，背面和侧面也必须共同解释同一个对象。单幅画面自然，只完成了其中一部分。

这句请求的困难来自两类关系。一类在单个输出内部：视频的各帧要描述同一个持续存在的对象，三维内容的各个视角要描述同一个表面。另一类在不同输出之间：声音要落在画面事件上，说出的回答要与看到的证据一致。读完本章，读者应能说明每种成果需要哪些条件、模型怎样把这些条件接入生成计算、怎样把生成结果解码成可交付的媒体，以及怎样核验上述两类关系。

要回答这些问题，先得知道“做得好不好”怎样度量，以及难点具体落在哪里。本章第一节把开篇请求拆成几项可以度量的挑战，并说明经典方法和学习方法分别能处理到哪一步；之后各节按“哪些内容已经固定、哪些内容由模型生成”逐步改变设定，前一节固定的内容在后一节变成可变量，这就是各节相连的原因。阅读本章需要第一部分的视频运动与时间连续性度量、多模态对齐与融合以及相机几何，也需要\ref{chap:vision-multimodal-models}章的生成模型接口。

生成过程沿用\ref{sec:vision-models-media-generation}节的潜表示、条件连接和采样接口，流匹配的速度目标见\ref{sec:vision-models-flow-matching}节，扩散模型的前向加噪与反向去噪可参考\BookChapterRef{chap:pgm-generative}{卷六“生成模型”章}。下文用$s\in[0,1]$表示流匹配的生成路径参数，用$\ell$表示扩散模型的噪声水平，用$t$或$\tau$表示媒体中的时间，避免把“第几次去噪”与“视频第几秒”混在一起。单幅图像的创建、编辑与观测场景的新视角合成见\ref{chap:vision-image-generation}章；本章增加跨时间、跨输出和跨视角的约束。

统一模型也要按实际输出入口使用。例如\ref{sec:vision-models-showo2}节的Show-o2可以在共享主干中形成文字与视觉输出，但视频能力取决于具体变体的训练范围，也没有因此获得音频输出。共享主干和多模型组合都要接受同样的成果核验，本章因此按成果而不是按主干组织。

\section{多模态生成的问题与挑战}
\label{sec:vision-multigen-challenges}

多模态生成的输入是创作条件，输出是一件或几件媒体成果。输入可以是文字；可以再加首帧、参考图或几帧已知画面；可以是待修改的原片段及修改指令或掩码；可以是一段已经固定、等待配声的视频；也可以是一组带相机参数的观测，或者用户对一段视频提出的问题。输出相应地是帧序列、与视频等长的声轨波形、带声轨的视频、可以从任意相机渲染的三维表示（辐射场、网格或高斯集合），以及文字与语音形式的回答。与\ref{chap:vision-multimodal-understanding}章的理解任务相比，这里最大的不同是同一请求有许多合法输出：“玩具狗绕花瓶走一圈”可以从左边绕，也可以从右边绕，可以走得快些或慢些。评价因此很少能把输出与唯一的参考逐像素比较，常用的度量分成三类：把一批生成结果与一批真实媒体在特征分布上比较，把每个结果与它的输入条件比较，以及检查一件成果内部或几件成果之间的关系。下面先给出本章用到的主要指标，再逐项讨论挑战。

视频的分布级指标中最常用的是\term{Fréchet视频距离}（Fréchet video distance, FVD）。单个生成视频没有唯一的正确答案，但一个好的生成模型产生的一批视频，其总体统计应当接近同类真实视频。FVD先把每段视频送入在Kinetics-400动作识别数据上训练的I3D网络（结构见\ref{sec:vision-video-i3d}节），取末层输出作为该视频的特征向量；再把真实视频的特征和生成视频的特征各自拟合成多元高斯分布，均值和协方差分别为$(\bm\mu_{\mathrm r},\bm\Sigma_{\mathrm r})$与$(\bm\mu_{\mathrm g},\bm\Sigma_{\mathrm g})$；最后计算两个高斯之间的Fréchet距离\scite{vision-fvd2018}计算式为
\begin{equation}
  \mathrm{FVD}
  =\lVert\bm\mu_{\mathrm r}-\bm\mu_{\mathrm g}\rVert_2^2
   +\operatorname{Tr}\!\left(
     \bm\Sigma_{\mathrm r}+\bm\Sigma_{\mathrm g}
     -2(\bm\Sigma_{\mathrm r}\bm\Sigma_{\mathrm g})^{1/2}
    \right)\text{。}
  \label{eq:vision-multigen-fvd}
\end{equation}
这个式子是两个多元高斯之间2-Wasserstein距离的平方。2阶Wasserstein距离的平方是把一个分布“搬运”成另一个分布所需的最小平均平方位移；其定义及多元高斯的闭式可参考\BookChapterRef{chap:information-theory-and-statistical-geometry}{卷2《数学准备：概率、因果与决策》的“概率分布的分析与优化”章}。本章把这一闭式用于视频特征。第一项衡量两批特征的中心偏了多远；第二项衡量它们的散布形状差多少。矩阵$\bm\Sigma_{\mathrm r}\bm\Sigma_{\mathrm g}$一般不对称，它的平方根按$\operatorname{Tr}[(\bm\Sigma_{\mathrm r}\bm\Sigma_{\mathrm g})^{1/2}]=\operatorname{Tr}[(\bm\Sigma_{\mathrm r}^{1/2}\bm\Sigma_{\mathrm g}\bm\Sigma_{\mathrm r}^{1/2})^{1/2}]$理解，右侧是对称半正定矩阵的平方根，迹为非负实数。当两个协方差都是对角阵时，计算可以化到逐维：设第$k$维的标准差为$\sigma_{\mathrm r,k}$和$\sigma_{\mathrm g,k}$，则$(\bm\Sigma_{\mathrm r}\bm\Sigma_{\mathrm g})^{1/2}$的第$k$个对角元为$\sigma_{\mathrm r,k}\sigma_{\mathrm g,k}$，迹项变为$\sum_k(\sigma_{\mathrm r,k}^2+\sigma_{\mathrm g,k}^2-2\sigma_{\mathrm r,k}\sigma_{\mathrm g,k})=\sum_k(\sigma_{\mathrm r,k}-\sigma_{\mathrm g,k})^2$。

用一个二维特征的教学构造检查各项含义。设真实特征均值为$(0,0)$、两维标准差都是1；生成特征均值为$(0.3,0.4)$、标准差为$(0.5,1)$，协方差都取对角。均值项为$0.3^2+0.4^2=0.25$，迹项为$(1-0.5)^2+(1-1)^2=0.25$，FVD为0.5。第二部分的0.25来自生成视频在第一维上散布过窄，这对应“生成结果彼此太像”的多样性不足；即使把均值完全对准，这一项也不会消失。FVD越低越好，但它有几条使用边界。多元高斯只是特征分布的近似；协方差要从样本估计，样本少时估计误差可能大于模型之间的差别；所用网络版本、帧数、分辨率和样本数改变，数值就不可比较。它还是整批视频的单一数字，不能指出哪一段视频在哪一帧出了问题。

为了定位失败，\ref{sec:vision-multigen-vbench}节的VBench把视频质量拆成16个维度，每个维度用一种专门工具给出逐视频分数。以主体一致性为例，设$\bm d_t$是第$t$帧的DINO特征并已归一化为单位长度，$T$为帧数，分数为
\begin{equation}
  S_{\mathrm{subject}}
  =\frac{1}{T-1}\sum_{t=2}^{T}
    \frac12\left(\langle\bm d_1,\bm d_t\rangle+\langle\bm d_{t-1},\bm d_t\rangle\right)\text{。}
  \label{eq:vision-multigen-subject-consistency}
\end{equation}
与首帧比较约束长时间的外观保持，与前一帧比较约束相邻帧的连续变化，两项平均\scite{vision-vbench2024}。三帧教学构造中，若$\langle\bm d_1,\bm d_2\rangle=0.9$、$\langle\bm d_1,\bm d_3\rangle=0.7$、$\langle\bm d_2,\bm d_3\rangle=0.8$，第2帧的项为$(0.9+0.9)/2=0.9$，第3帧为$(0.7+0.8)/2=0.75$，分数为0.825。一段几乎不动的视频也能得到接近1的分数，所以这一维必须与动态程度一起看。

声音的分布级指标沿用同一思路。\term{Fréchet音频距离}（Fréchet audio distance, FAD）把音频切成1秒的窗，每0.5秒取一窗，送入VGGish音频分类网络，取分类层之前的128维激活作为嵌入，再对生成音频与参考音频的嵌入按式\eqref{eq:vision-multigen-fvd}计算Fréchet距离\scite{vision-fad2019}。VGGish的输入是由1秒音频形成的64维对数梅尔特征序列，相位信息不进入嵌入，因此只改变相位的失真可能检测不到。MMAudio等后续工作还用PaSST、PANNs等音频网络的嵌入计算同一距离，统称Fréchet距离（FD），不同嵌入的数值不能互相比较\scite{vision-mmaudio2025}。分布距离不看配对关系：把所有声轨随机交换给别的视频，FAD不变。检查“这段声音是否属于这段视频”需要逐片段的指标。一种做法是用预训练音频分类器把生成音频和该视频的原始声轨分别变成类别概率分布$\bm p_i^{\mathrm{gen}}$和$\bm p_i^{\mathrm{ref}}$，再对$n$个测试片段平均二者的KL散度，例如$\frac1n\sum_i\operatorname{KL}(\bm p_i^{\mathrm{ref}}\|\bm p_i^{\mathrm{gen}})$；散度的方向和概率的归一化方式随所沿用的实现而定，比较模型时须保持一致。另一种做法不需要原声轨：用\ref{sec:vision-models-imagebind}节的ImageBind分别编码视频与生成音频，取两个嵌入的余弦相似度，称为IB分数。两者都在语义层面判断“是不是这种声音”，不回答“是不是在这一刻响”。

音画同步需要时间上的指标。逐事件的检查用\ref{sec:vision-multigen-audio-check}节式\eqref{eq:vision-event-sync-error}的起声偏移：先把视频接触与音频起声配对，再计算偏移及其平均绝对值。逐事件标注成本高时，可以用自动估计的整段偏移：DeSync用在音画偏移检测上训练的Synchformer预测整段音频相对视频的时间偏移，取绝对值，单位为秒；八秒音频时，MMAudio论文对开头和结尾两段4.8秒的截取分别预测再平均。DeSync越小越同步，但它只给出整段的一个偏移，不能指出哪次碰撞出错。

三维内容的指标取决于有没有真实形状。没有参考形状时，只能检查渲染结果是否符合文字。DreamFusion使用CLIP R-Precision：从固定的仰角和若干方位角渲染生成对象，用CLIP判断渲染图在一组文字描述中最匹配哪一条，统计取到正确描述的比例\scite{vision-dreamfusion2023}。设共有$N$条描述$c_1,\dots,c_N$，第$i$个对象的渲染图为$\bm x_i$，图像与文字编码器为$f_{\mathrm{img}},f_{\mathrm{txt}}$，则
\begin{equation}
  \mathrm{R@1}
  =\frac1N\sum_{i=1}^N
   \mathbb 1\!\left[
    \argmax_{j}\cos\bigl(f_{\mathrm{img}}(\bm x_i),f_{\mathrm{txt}}(c_j)\bigr)=i
   \right]\text{。}
  \label{eq:vision-multigen-clip-rprecision}
\end{equation}
论文同时在带颜色的渲染图和去掉纹理的几何渲染图上计算这一比例，因为把纹理画在平板几何上，颜色渲染也可能得到高分。有参考形状时可以直接比较几何。Wonder3D在Google Scanned Objects中的30个对象上，用倒角距离与体积交并比比较重建表面和真实表面，用PSNR、SSIM与LPIPS比较生成的颜色视图与真实视图\scite{vision-wonder3d2024}后三者的定义见\ref{sec:vision-restoration-quality}节。设$P,Q$是从两个表面上采样的点集，\term{倒角距离}（Chamfer distance）常写作
\begin{equation}
  d_{\mathrm{CD}}(P,Q)
  =\frac{1}{|P|}\sum_{\bm p\in P}\min_{\bm q\in Q}\lVert\bm p-\bm q\rVert_2^2
   +\frac{1}{|Q|}\sum_{\bm q\in Q}\min_{\bm p\in P}\lVert\bm q-\bm p\rVert_2^2,
  \label{eq:vision-multigen-chamfer}
\end{equation}
两项分别检查一个表面上的点能否在另一个表面附近找到对应点，二者合起来同时惩罚多出的部件和缺失的部件。教学构造中取$P=\{(0,0),(1,0)\}$、$Q=\{(0,0),(1,0.2)\}$，两个方向各有一个点的最近平方距离为0.04，平均后各为0.02，倒角距离为0.04。不同实现可能用距离而不是平方距离，或先把形状缩放到单位尺度，数值只在同一实现内可比。\term{体积交并比}（volume intersection over union, volume IoU）把两个形状体素化后计算交集体素数与并集体素数之比，1表示完全重合。生成的多个视角之间是否互相一致，目前没有一个被普遍采用的单一指标；常用的检查是在已知相机下把一个视角中的点按深度投影到另一个视角，比较对应位置的颜色或法线，或者干脆把多视角图重建成三维再看重建误差。

表\ref{tab:vision-multigen-metrics}把这些指标按比较对象归在一起。表中每一行回答的问题不同，一个指标的高分不能代替另一行。

\begin{table*}[tbp]
 \centering
 \small
 \begin{tabularx}{\linewidth}{@{}p{20mm}p{30mm}p{34mm}X@{}}
  \toprule
  成果 & 指标 & 比较对象 & 不能说明的内容 \\
  \midrule
  视频 & FVD & 生成与真实视频的I3D特征分布 & 单个视频的哪一帧失败；是否符合文字 \\
  视频 & VBench各维度 & 每段视频的帧间特征、光流、检测结果等 & 事件是否合理；不同维度之间的取舍 \\
  音频 & FAD或FD & 生成与参考音频的嵌入分布 & 声音是否属于该视频；何时发声 \\
  音频 & KL、IB分数 & 与原声轨或视频的逐片段语义 & 起声时刻 \\
  音画 & 起声偏移、DeSync & 事件时间或整段偏移 & 声源是否正确；DeSync不能定位到事件 \\
  三维 & CLIP R-Precision & 渲染图与文字描述 & 几何是否正确；背面是否合理 \\
  三维 & 倒角距离、体积交并比 & 重建表面与真实表面 & 纹理质量；需要真实形状 \\
  \bottomrule
 \end{tabularx}
 \caption{本章使用的主要评价指标。分布级指标针对一批成果，逐片段指标针对单件成果；各指标依赖的特征网络或工具改变时，数值不可比较。}
 \label{tab:vision-multigen-metrics}
\end{table*}
\FloatBarrier

有了这些度量，开篇请求的难点可以拆成五项。每一项都先在玩具狗与花瓶的例子中出现，然后看经典方法能处理到哪一步，学习方法改变了什么。

\subsection{帧间依赖带来的时间一致性困难}
\label{sec:vision-multigen-challenge-temporal}

一段五秒、每秒16帧的视频有80幅画面。每一幅单独看都可能是合理的玩具狗图像，但相邻两幅还必须是同一个连续世界的两个瞬间：背上的花纹不能逐帧换位置，腿的摆动要有一致的节奏，玩具狗贴近花瓶时两者不能互相穿透。逐帧独立生成会同时破坏这三点，常见表现是画面闪烁、外观漂移和运动抖动。运动是否合理更难判断：步幅忽大忽小、碰到花瓶却没有任何反应，都可能出现在每帧清晰的视频里。

经典视频处理有两类工具与此相关。一类是运动估计：\ref{sec:vision-videofound-motion}节的块匹配与光流给出相邻帧之间每个位置移动到哪里。另一类是\ref{sec:vision-videofound-consistency}节的时间连续性度量：闪烁度量比较静止区域相邻帧的亮度变化；扭曲误差先用光流把前一帧变形到后一帧的位置，再在非遮挡区域比较二者的差异，差异大说明外观没有沿运动保持。在复原或增强已有视频时，这些工具可以直接构成约束，例如沿运动轨迹做时间滤波。生成任务却没有原视频，运动本身就是要生成的内容；光流只能在生成之后用来检查，不能告诉模型玩具狗应当怎样走。

学习方法把“整段视频由一个模型共同决定”作为出发点。\ref{sec:vision-multigen-create-video}节的Lumiere在网络内部同时下采样时间和空间，使远隔的帧在最粗层级共同参与计算；Wan先用视频自编码器压缩相邻帧的冗余，再让注意力覆盖全部潜位置。运动先验来自大规模视频数据，而不是手工的平滑假设。经典度量并没有被淘汰，而是以学习工具的形式回到评价中：VBench的运动平滑度借插帧模型检查帧间运动，动态程度借RAFT光流衡量运动幅度。视频修改面对同一困难的另一面：\ref{sec:vision-multigen-edit-video}节要求新内容逐帧一致，未修改的部分还要保持原样。运动在物理上是否合理，最终还要交给\ref{sec:vision-multigen-vbench2}节按事件提问的评价。

\subsection{模态时间尺度差异带来的同步困难}
\label{sec:vision-multigen-challenge-sync}

玩具狗在1.0秒碰到花瓶，碰撞声就应在1.0秒附近开始。困难在于两种媒体的时间单位相差很大。每秒24帧的视频，一帧覆盖约41.7毫秒；每秒8帧的语义特征，一个位置覆盖125毫秒；16 kHz的音频每秒有16000个采样点，压缩到潜空间后每秒仍有约31个位置。若生成网络把视频第$k$个位置与音频第$k$个位置当作同一时刻，就会产生成倍的时间错位。更细的困难是事件粒度：一次碰撞在画面里可能只占一两帧，平均到整段的特征里几乎看不出来。

经典的多模态对齐处理的是两条已经存在的序列。\ref{sec:vision-alignment-dtw}节的动态时间规整在两条序列之间寻找一条单调的对应路径，使累积的逐点代价最小，允许一方局部快一些、另一方局部慢一些；\ref{sec:vision-alignment-alignment}节还区分了对齐的粒度，即在帧、片段还是整段层面建立对应。在配声的生成任务中，这两点各自起作用。粒度的区分直接决定评价方式：DeSync给的是整段粒度的偏移，起声偏移是事件粒度的检查。动态时间规整则需要两条序列都已存在，而生成时至少有一条尚未产生，它不能把同步“写进”生成过程；它仍可在核验时用来匹配事件，但它的单调路径只能拉伸时间，不能表达多出或漏掉的一次碰撞，这些要另行计数。

学习方法把时间对应放进生成的条件接入中。MMAudio按比例缩放视觉序列的旋转位置编码，并加入每秒24个位置的同步特征逐位置调制音频表示；Movie Gen Audio让每个音频潜位置读取时间上最近的视频帧嵌入；Ovi和LTX-2在两个分支的每一层用带时间位置的交叉注意力双向交换信息，这是\ref{sec:vision-alignment-intermediate-fusion}节中间融合在生成模型中的形态。位置编码只让“同一时刻”可以被表达，正确的对应仍要从配对数据中学到。画面固定时的配声由\ref{sec:vision-multigen-audio}节回答，画面与声音都可变时由\ref{sec:vision-multigen-joint}节回答，同步的核验在\ref{sec:vision-multigen-audio-check}节。

\subsection{生成长度增长带来的状态保持困难}
\label{sec:vision-multigen-challenge-long}

本章的视频模型一次覆盖约五秒，音频模型一次最多约30秒。一个一分钟的玩具狗故事因此要由十几段视频接成。每段内部可以很连贯，段与段之间却可能出现三类问题：玩具狗被花瓶遮挡后重新出现时换了花纹，即身份不保持；花瓶的位置每段漂移几个像素，累积起来离开了原位，即场景不保持；背景环境声在段落交界处突然变化。时间越长，需要记住的状态越多，而模型一次能看到的范围是固定的。

经典视频分析为长视频提供了结构单位。\ref{sec:vision-videofound-structure}节的镜头边界检测把视频切成镜头，关键帧选择用少数画面代表一个镜头，背景建模把静止的场景与运动的前景分开；\ref{sec:vision-videofound-consistency}节的身份连续性度量检查同一对象在不同时间是否仍被识别为同一对象。这些方法描述已有视频的结构。在生成任务里，它们转而成为规划和检查的单位：先按镜头规划，再以关键帧作为每段的条件，最后用身份度量检查跨段一致。但经典方法本身不生成任何新的帧。

学习方法目前的主要手段是把已知内容作为条件接到下一段。\ref{sec:vision-multigen-extension}节用Wan的已知帧掩码续接视频；Movie Gen Audio把前一段末尾的音频作为下一段的上下文，或者在每个生成步对重叠窗口共同更新；Lumiere的空间超分辨率也用重叠窗口协调相邻片段。这些做法改善了交界处的衔接，但都没有保存一个持久的场景状态，误差会沿续接链条传递。这项挑战在本章只得到部分回答，本章小结会回到它。

\subsection{视角变化带来的几何一致性困难}
\label{sec:vision-multigen-challenge-geometry}

把玩具狗做成三维资产，要求从正面、侧面和背面看到的画面都来自同一个表面。二维图像生成模型分别画这三幅图，常常各自合理却互相矛盾：耳朵形状不同，或者背面也长出一张脸。只有一幅正面照片时，背面根本没有被观测；玩具狗在运动时，不同时刻拍到的腿部位置又不能拟合成同一个静止表面。

经典的多视图几何在有观测的条件下给出严格的一致性。\ref{sec:vision-formation-two-view}节的对极几何说明，同一个三维点在两幅图中的投影必须满足对极约束，已知两个相机时可以三角化出它的位置；\ref{sec:vision-formation-multi-view}节的运动恢复结构与光束法平差把这一点推广到多幅图，同时估计相机和三维点。\ref{sec:vision-generation-novel-view}节的NeRF与3D Gaussian Splatting则用可微渲染把一组带位姿的照片拟合成场景表示，表示本身是共享的，所以各视角自动一致。这些方法的前提是观测覆盖了要重建的表面，并且图像之间能找到对应；它们不会为没有拍到的背面创作内容，也不能从一句文字出发。

学习方法把二维生成先验和三维一致性结合起来，结合的位置有三种。DreamFusion保留一个共享的三维表示，用冻结的二维扩散模型为每幅渲染图提供修改方向；Wonder3D让二维扩散模型一次生成多个视角，在生成时就用跨视角注意力协调；4D Gaussian Splatting在有观测时不再需要生成先验，用形变场把多时刻观测连成一个动态场景。三者分别对应\ref{sec:vision-multigen-3d}节的三种观测条件。

\subsection{迭代采样带来的制作预算困难}
\label{sec:vision-multigen-challenge-budget}

本章的生成模型都通过多步迭代得到结果。Wan的报告通常采样约50步，若每步做无分类器引导，还要多算一次去条件预测，一段视频就需要约100次主干前向；LTX-2分别调节文字与模态引导时每步三次前向；DreamFusion为一条文字请求优化15000步，论文在4块TPUv4芯片上用约1.5小时。单次成本之外，还有候选数量：一条请求往往要生成多个候选，再挑选、局部重做。只比较最终最好的样本，会把这些成本都藏起来。

经典方法降低计算的主要思路是由粗到细与去除冗余。\ref{sec:vision-transforms-pyramid}节的图像金字塔先在低分辨率上求解全局结构，再逐级补充细节；视频编码利用相邻帧的相似性只存储变化部分。这两种思路都在学习方法中延续：Lumiere先在$128\times128$生成整段运动，再用空间超分辨率补细节；Wan的视频自编码器把时间压缩4倍、空间各压缩8倍，生成网络面对的位置数降到原来的约$1/256$，再在压缩后的潜网格上迭代；LTX-2还在基础潜变量生成后再做潜空间的空间上采样与分块细化。压缩减少了每步的计算，却引入了自编码器的重建误差；级联减少了基础模型的计算，却把协调问题留给了超分辨率阶段。

成本因此不是附带信息，而是比较模型时必须固定的条件。各节在介绍模型时给出论文报告的步数、耗时与规模，\ref{sec:vision-multigen-budget}节说明怎样在同一预算口径下比较制作流程。

这五项挑战决定了本章的展开顺序。时间一致性在只有画面的设定中最纯粹，所以\ref{sec:vision-multigen-video}节先处理视频的创建、修改和延伸，延伸部分同时触及长时状态保持。画面固定之后，同步困难只剩一个方向，由\ref{sec:vision-multigen-audio}节在配声中回答；画面和声音都可变时，同步要求两路在生成中互相读取，这是\ref{sec:vision-multigen-joint}节的问题。几何一致性把一致性的对象从时间轴换到视角，由\ref{sec:vision-multigen-3d}节按观测多少依次讨论。\ref{sec:vision-multigen-response}节的语音回答把同步要求推到播放时刻：说出的内容要等证据到达，又不能拖得太久。最后，\ref{sec:vision-multigen-evaluation}节用本节的指标和预算口径汇总验收，并说明哪些挑战仍未解决。

\section{视频生成与编辑}
\label{sec:vision-multigen-video}

视频生成与编辑的输出是一段帧序列，输入则随任务不同：创建视频时只有文字，或再加一幅首帧；修改视频时还有原片段，以及说明改哪里的指令或掩码；延伸视频时有一段已知帧，待生成的是它之前、之后或之间的时间范围。三类任务共同的难点是时间一致性：同一只玩具狗要在每一帧保持外观，运动要连贯，而一段五秒、每秒16帧的视频就有80幅画面，网络很难在显存允许的范围内同时看到全部像素。修改和延伸还多一个难点：模型必须知道哪些内容属于输入、应当保留，哪些内容允许改变。

本节按任务类型划分小节。创建视频时，两条路线对“怎样让整段视频共同参与计算”给出不同回答：Lumiere在像素空间的网络内部压缩时间，Wan先用视频自编码器把整段视频压成潜网格，再在潜网格上用Transformer处理全部位置。修改视频时，两条路线对“怎样告诉模型改哪里”给出不同回答：Movie Gen Edit从文字指令中学习修改关系，Lumiere视频修补用时空掩码直接指定区域。延伸与补全使用Wan的已知帧掩码接口，它把“哪些帧已知”编码为条件，与修补中“哪些像素已知”的做法相对应。

\subsection{视频的创建}
\label{sec:vision-multigen-create-video}

创建视频的输入是文字，有时还有一幅首帧；输出是规定时长、帧率和画幅的帧序列。文字规定场景和事件，首帧还规定故事开始时的外观与布局。若请求包含精确轨迹，则需要模型接收随时间变化的位置或运动条件；仅支持文字和首帧的模型没有获得这条数值轨迹。相同的“镜头绕行”，也可能由相机移动、对象旋转或两者共同运动形成，不能只检查画面里出现了转动。

开始生成前，制作记录要把时长换算成帧数。一个教学记录规定四秒、每秒24帧、固定画幅，以及“玩具狗从花瓶左侧走到右侧，花瓶保持不动”。按每帧占据一个播放时间间隔的约定，需要96帧，第$i$帧的起始时刻为$(i-1)/24$秒。若只规定首尾采样时刻相距四秒，则包含两端会得到97帧。模型自身还可能要求帧数符合特定分组规则，制作记录应把采样范围、容器播放时长和模型输入长度分别写清。

两个代表模型都要解决同一个矛盾：时间一致性要求网络同时看到整段视频，而计算预算不允许在全分辨率下这样做。Lumiere保留像素表示，在网络内部同时对空间和时间做下采样，在最粗的时空网格上处理整段运动，再用分窗口的空间超分辨率补细节。Wan把压缩前移到独立训练的视频自编码器中，生成网络只面对压缩后的潜网格，并用流匹配训练。前者的代价落在级联的超分辨率阶段，后者的代价落在自编码器的重建误差和潜空间序列长度上。

\subsubsection{Lumiere}
\label{sec:vision-multigen-lumiere}

Lumiere之前的文字生成视频系统常采用级联：基础模型先生成时间上稀疏的关键帧，再由时间超分辨率模型在关键帧之间插入画面。这种做法在关键帧间距较大时会遇到具体困难。设玩具狗的腿每0.25秒摆动一个周期，关键帧若每0.5秒取一幅，相邻两幅中腿的位置几乎相同；插帧模型只看到两幅相似画面，无法判断其间是静止、摆了一次还是摆了两次，生成的步态因此可能变慢或抖动。这是时间采样不足造成的混叠，补帧阶段缺少决定运动所需的信息\scite{vision-lumiere2024}

Lumiere的主要想法是让基础模型一次处理整段视频的全部帧，在网络内部对时间维也做下采样和上采样。时间被压短后，远隔帧的信息能在较小表示上共同参与计算，整段运动在最粗的时空网格上就已决定；随后逐层上采样，恢复完整帧率。由于基础模型直接输出完整帧率序列，后续阶段只需补空间细节，不再需要时间插帧，上述混叠来源就被移除了。作者把这种网络称为\term{时空U-Net}（space-time U-Net, STUNet），即在空间和时间两个方向上同时做多尺度下采样与上采样的U形去噪网络。

模型从预训练的Imagen文字生成图像模型出发。Imagen在像素空间工作，由一个低分辨率基础模型和空间超分辨率级联组成。Lumiere基础模型的输入是含噪视频$\bm X_\ell\in\mathbb R^{80\times128\times128\times3}$、噪声水平$\ell$和文字表示，输出同形状的噪声预测；80帧按每秒16帧对应五秒。论文给出的结构信息有三条。第一，时间下采样和上采样模块紧跟在预训练网络的每个空间缩放模块之后，所以空间每缩小一级，时间也同时缩短。第二，除最粗层级外，每个层级插入分解的时空卷积，即先做空间卷积、再沿时间做一维卷积；分解的好处是新增参数只有一维时间核，原有空间卷积可以原样保留。第三，最粗层级插入时间注意力，让每个空间位置在全部剩余时间位置之间交换信息；注意力的成本随序列长度平方增长，只有在时间已被压短的最粗层级才负担得起。新增模块以残差方式接入，其输出投影初始化为零映射。论文没有给出层级数、每一级的时间下采样倍数和参数量。为了看清张量形状怎样变化，以每个层级时空各减半作为教学假设，形状依次为$80\times128^2\to40\times64^2\to20\times32^2\to10\times16^2$，最粗层级只剩10个时间位置；表\ref{tab:vision-multigen-lumiere-stunet}按这个假设列出各层级的组成，其中形状一列不是论文发布的配置。

\begin{table*}[tbp]
 \centering
 \small
 \begin{tabularx}{\linewidth}{@{}p{24mm}p{30mm}p{34mm}p{26mm}X@{}}
  \toprule
  层级组 & 预训练空间部分 & 新增时间层 & 时空采样 & 输出形状（教学假设） \\
  \midrule
  编码第1级 & 空间卷积块，冻结 & 分解时空卷积 & 空间、时间各下采样 & $40\times64\times64$ \\
  编码第2级 & 空间卷积块，冻结 & 分解时空卷积 & 空间、时间各下采样 & $20\times32\times32$ \\
  编码第3级 & 空间卷积块，冻结 & 分解时空卷积 & 空间、时间各下采样 & $10\times16\times16$ \\
  最粗层级 & 空间卷积与注意力，冻结 & 时间注意力 & 无 & $10\times16\times16$ \\
  解码各级 & 空间卷积块，冻结，接收跳跃连接 & 分解时空卷积 & 空间、时间各上采样 & 逐级恢复到$80\times128\times128$ \\
  \bottomrule
 \end{tabularx}
 \caption{Lumiere时空U-Net的层级组成。时间层的种类与位置依据论文；论文未给出层级数、每级下采样倍数与参数量，表中层级数和输出形状是便于手算的教学假设，通道数随原图像U-Net而定，论文未单独列出。}
 \label{tab:vision-multigen-lumiere-stunet}
\end{table*}
\FloatBarrier

训练使用约3000万段带文字描述的视频，每段80帧、每秒16帧。原图像模型的权重保持固定，只更新新增的时间参数。论文没有单列改写后的目标函数；按其沿用的扩散概率模型训练，可用常见的噪声预测形式写成
\begin{equation}
  L_{\mathrm{Lum}}
  =\mathbb E_{\bm X,c,\ell,\bm\epsilon}
   \left[
    \left\lVert
     \widehat{\bm\epsilon}_{\bm\theta}(\alpha_\ell\bm X+\sigma_\ell\bm\epsilon,\ell,c)-\bm\epsilon
    \right\rVert_2^2
   \right],
  \label{eq:vision-lumiere-objective}
\end{equation}
其中$\bm X$是训练视频，$c$是文字条件，$\bm\epsilon$是与视频同形状的标准高斯噪声，$\alpha_\ell,\sigma_\ell$由噪声日程给出，$\bm\theta$只包含新增时间参数。冻结图像权重保留了大规模图文训练得到的外观先验，新增参数只需学习时间关系。时间下采样会混合或选择帧，网络初始化不能简单解释为“每帧恰好等于原图像模型”；原方法把时间采样模块初始化为最近邻下采样和上采样，再加上零初始化的输出投影，使训练开始时的网络行为接近逐帧运行图像模型。

推理时从整段噪声视频出发，反复运行时空网络并执行采样更新。一次前向能看完整段，不表示一次前向已经完成采样。基础模型输出$80\times128\times128$的视频后，空间超分辨率模型把每帧提升到$1024\times1024$。高分辨率网络受显存限制，只能处理较短时间窗口。若各窗口独立完成再拼接，边界两侧可能出现不同的毛色或光照。Lumiere采用MultiDiffusion的做法，在每一步采样中处理重叠窗口，再对重复覆盖的位置融合预测，使下一步各窗口从同一段全局状态继续采样；原文的时间窗口重叠两帧，窗口长度论文未给出。

设窗口$i$选择全局视频中的若干连续帧，选择算子为$\bm P_i$；把当前一步的窗口预测展平成向量$\bm y_i$。希望新的全局状态$\bm z$在每个窗口内都接近相应预测，可以求解
\begin{equation}
  \min_{\bm z}\sum_i\lVert\bm P_i\bm z-\bm y_i\rVert_2^2,
  \qquad
  \widehat{\bm z}
  =\left(\sum_i\bm P_i^\top\bm P_i\right)^{-1}
     \sum_i\bm P_i^\top\bm y_i\text{。}
  \label{eq:vision-window-consensus}
\end{equation}
这里假设所有全局位置至少被一个窗口覆盖，且$\bm P_i$只负责取出对应坐标。对$\bm z$求导并令梯度为零，就得到右侧的正规方程。矩阵$\sum_i\bm P_i^\top\bm P_i$是对角矩阵，每个对角元素等于该位置被覆盖的次数，因此解就是对重叠预测逐位置求平均。

例如八帧教学序列采用窗口1至5和4至8，第4、5帧被预测两次。若第4帧某个颜色分量的两次预测为0.4和0.6，合并值为0.5；只被覆盖一次的分量直接保留。图\ref{fig:vision-window-consensus}展示了覆盖计数的来源。这个平均发生在每步采样的状态协调中，不能用“最终把两段成片淡入淡出”代替。

把这一步写成算法时，还可以给每个窗口内的位置一个非负权重$w_{i,k}$，等权平均是全部权重取1的特例；\ref{sec:vision-multigen-movie-audio}节的Movie Gen Audio就使用中间高、两端低的三角形权重。若把式\eqref{eq:vision-window-consensus}中的平方误差按$w_{i,k}$加权，同样令梯度为零，每个全局位置的解就是覆盖它的各个预测按权重的加权平均，算法\ref{alg:vision-multigen-window-consensus}给出完整步骤。

\begin{algorithm}[重叠窗口的逐步共同更新]
  \label{alg:vision-multigen-window-consensus}
  \begin{algorithmic}[1]
    \Require 全局含噪状态$\bm z^{(0)}$；覆盖全部位置的窗口集合$\{\bm P_i\}$及窗口内权重$w_{i,k}\geq0$；单窗口去噪网络；采样步数$K$
    \Ensure 全局生成结果$\bm z^{(K)}$
    \For{$k=0,1,\dots,K-1$}
      \For{每个窗口$i$}
        \State 取出$\bm P_i\bm z^{(k)}$，运行网络与一步采样更新，得到窗口预测$\bm y_i$
      \EndFor
      \For{每个全局位置$u$}
        \State 累加分子$a_u\gets\sum_{(i,j):\,\bm P_i\text{的第}j\text{个坐标为}u}w_{i,j}\,y_{i,j}$
        \State 累加分母$b_u\gets\sum_{(i,j):\,\bm P_i\text{的第}j\text{个坐标为}u}w_{i,j}$
        \State 若$b_u=0$则报告窗口布置错误；否则$z^{(k+1)}_u\gets a_u/b_u$
      \EndFor
    \EndFor
  \end{algorithmic}
\end{algorithm}

用上面八帧的例子复核算法。若两个五帧窗口都采用三角形权重$(1,2,3,2,1)$，第4帧在窗口一中是第4个位置，权重为2，在窗口二中是第1个位置，权重为1。两个预测仍为0.4和0.6时，合并值为$(2\times0.4+1\times0.6)/(2+1)\approx0.467$，比等权平均更靠近权重较大的窗口一。权重的作用是让窗口边缘的预测少占比例，因为边缘位置只看到单侧的上下文。

\begin{figure}[htbp]
  \centering
  \begin{tikzpicture}[x=1cm,y=1cm,font=\figurefont\small,
  cell/.style={draw=BookInk,line width=0.7pt,minimum width=7mm,
    minimum height=7mm,inner sep=0pt},
  arr/.style={-{Stealth[length=2mm]},BookInk,line width=0.8pt}]
  \node[anchor=east] at (1.65,2.45) {帧位置};
  \node[anchor=east] at (1.65,1.6) {窗口一};
  \node[anchor=east] at (1.65,0.65) {窗口二};
  \node[anchor=east] at (1.65,-0.45) {覆盖次数};
  \foreach \i in {1,...,8} {
    \pgfmathsetmacro{\xx}{1.5+0.95*\i}
    \node at (\xx,2.45) {\i};
    \ifnum\i<6
      \node[cell,fill=FigureInput!15] at (\xx,1.6) {};
    \fi
    \ifnum\i>3
      \node[cell,fill=FigureOutput!15] at (\xx,0.65) {};
    \fi
    \ifnum\i=4
      \node at (\xx,-0.45) {2};
    \else
      \ifnum\i=5
        \node at (\xx,-0.45) {2};
      \else
        \node at (\xx,-0.45) {1};
      \fi
    \fi
  }
  \draw[BookAmber,dashed,line width=0.8pt] (4.85,0.15) rectangle (6.7,2.02);
  \draw[arr] (5.3,0.1) -- (5.3,-0.1);
  \node[anchor=west] at (1.65,-1.3) {同一位置：$(0.4+0.6)/2=0.5$};
\end{tikzpicture}

  \caption{重叠窗口的教学构造。两个五帧窗口在第4、5帧重叠，覆盖次数决定式\eqref{eq:vision-window-consensus}的逐位置平均。图中帧数用于手算，不是Lumiere的发布配置；原文空间超分辨率窗口在时间上重叠两帧。}
  \label{fig:vision-window-consensus}
\end{figure}
\FloatBarrier

让玩具狗绕花瓶走一圈时，可分别检查起终点是否连接、步态是否连贯、遮挡后外观是否保持。重叠融合能协调窗口边界，却不能证明模型理解了接触或三维形状。80帧是该设计所覆盖的时间范围，不能由此推出任意长度的叙事都能一次处理；短片段内部连贯与多镜头之间的身份保持，仍是不同的验收对象。

\Needspace{60mm}
\subsubsection{Wan}
\label{sec:vision-multigen-wan}

Wan的视频自编码器、扩散Transformer主干和条件接入方式已在\ref{sec:vision-models-wan}节说明，主干各规模的层数与宽度等配置也以该节为准，本节不再重复。这里只回顾使用时需要的结论：自编码器把视频在时间上压缩4倍、在高和宽上各压缩8倍，潜通道为16，第一帧只做空间压缩，作为时间起点单独保留；主干把潜网格切成时空块序列，用自注意力处理全部位置，文字经umT5编码后由交叉注意力接入；训练采用式\eqref{eq:vision-flow-loss}的流匹配速度目标\scite{vision-wan2025}。图像条件版本把首帧与其后补零的帧组成条件视频，经同一自编码器得到条件潜网格，再与含噪潜网格和重排后的已知位置掩码沿通道拼接；首帧还经CLIP图像编码器提取整体特征，通过独立的交叉注意力接入。

与Lumiere相比，Wan把时间一致性问题交给两处处理。自编码器先把冗余的相邻帧压缩，生成网络面对的位置数减少到原来的约$1/(4\times8\times8)$；主干的自注意力在压缩后的全部位置之间直接交换信息，不需要像STUNet那样在网络内部逐级缩短时间。因此Wan不需要级联的空间超分辨率，由自编码器解码直接得到目标分辨率的画面。报告给出1.3B和14B两种规模，训练使用每秒16帧、五秒长的片段，分辨率逐阶段提高到720p。

使用时首先组织输入。帧数要满足自编码器的分组约定，即第一帧之后的帧数是4的倍数。以$T=17$、$256\times256$画面为尺寸示例，目标潜网格为$5\times32\times32\times16$。文字生成只需要文字和目标尺寸；图像条件生成还需要首帧，首帧的宽高比应与目标画幅一致，否则缩放或裁剪会先改变参考内容。条件视频的其余帧填零，零表示未提供内容，不能把它当成“未来画面应为黑色”的监督目标。文字生成与图像条件生成使用不同的检查点，报告中图像编码器分支是在任务微调阶段才加入的，不能把文字生成权重直接配上首帧使用。

一般使用不再微调主干，主要调节采样设置。报告的推理通常使用约50步采样，并采用无分类器引导；引导强度越大，结果越贴近文字，但动作幅度和画面多样性可能下降。随机种子决定初始噪声，同一请求应保留多组种子的候选。采样结束后，潜网格由自编码器解码为帧序列；解码器的因果结构允许分块解码较长序列，但这不表示生成网络已经学会长时间状态保持。报告还指出，14B模型在单块高端GPU上未经优化的推理约需30分钟，这一成本应计入制作预算。

固定同一首帧，分别要求玩具狗转头、向前走和镜头环绕，可以观察首帧条件留下了什么自由度。转头要补出原来看不到的侧面，行走要决定腿部运动和接触，镜头环绕还要补背景与背面。首帧只给了起点，这些未知内容不会因“图生视频”这个任务名称而变成已知。常见失败包括：大幅运动时细节丢失，这也是报告列出的局限；转身后耳朵或花纹改变；文字与首帧冲突时，例如文字写“白色玩具狗”而首帧是棕色，模型可能在前几帧逐渐改色。冲突应在生成前消除，而不是期待模型自行选择。

两个模型可以在同一“玩具狗绕花瓶走一圈”的请求下比较，见表\ref{tab:vision-video-creation-compare}。Wan把时间压缩从生成网络移到自编码器，换来了单阶段的高分辨率输出和更多条件入口；代价是生成质量受自编码器重建误差限制，长序列的注意力成本随潜位置数的平方增长。Lumiere保留像素表示，冻结图像模型使训练只需学习时间参数，但五秒、$128\times128$的基础输出必须依靠分窗口超分辨率补细节，窗口之间只能靠逐步平均协调。

\begin{table*}[tbp]
 \centering
 \small
 \begin{tabularx}{\linewidth}{@{}p{24mm}XX@{}}
  \toprule
  比较项 & Lumiere & Wan \\
  \midrule
  生成空间 & 像素空间，基础模型$128\times128$，再级联空间超分辨率 & 视频自编码器的潜网格，时间压缩4倍、空间各8倍 \\
  整段时间的处理 & 网络内部时空下采样，最粗层级用时间注意力 & 自编码器先压缩，主干自注意力覆盖全部潜位置 \\
  训练目标 & 噪声预测扩散目标，冻结图像权重，只训练时间参数 & 流匹配速度回归，训练整个主干 \\
  条件入口 & 文字 & 文字；图像条件版本另有首帧潜网格、掩码和CLIP整体特征 \\
  输出范围 & 80帧、每秒16帧，超分辨率后$1024\times1024$ & 每秒16帧、五秒片段，最高720p \\
  主要代价 & 超分辨率窗口只能靠逐步平均协调 & 自编码器重建误差，注意力成本随潜位置数增长 \\
  \bottomrule
 \end{tabularx}
 \caption{视频创建两条路线的比较。各项依据两篇报告的说明，不是同一实验条件下的质量排名；两者都只覆盖约五秒的短片段。}
 \label{tab:vision-video-creation-compare}
\end{table*}
\FloatBarrier

\subsection{视频的修改}
\label{sec:vision-multigen-edit-video}

修改视频的输入是原片段和修改要求，输出是改动后的片段。难点有两个。第一，模型要区分哪些变化是任务要求，哪些变化应算误改。把花瓶换成金属杯主要改变对象外观；让玩具狗跳上桌子则改变动作、遮挡和接触，原背景中某些区域也会随之显露。保留范围必须根据请求判断，不能一律要求整段视频的所有像素不变。第二，真实的“原视频—编辑后视频”配对极少，模型很难直接从数据中学到修改关系。

两条路线对“怎样指定改哪里”给出不同回答。Movie Gen Edit只接收文字指令，由模型自己判断修改范围，监督则通过分阶段构造的训练配对逐步获得。Lumiere视频修补让用户提供时空掩码，直接规定允许改变的像素，训练配对可以从任意完整视频中遮去一部分自动产生。前者使用方便，但修改范围依赖模型理解；后者范围明确，但掩码必须随对象运动逐帧给出。

\Needspace{60mm}
\subsubsection{Movie Gen Edit}
\label{sec:vision-multigen-movie-edit}

只有“文字—视频”配对时，模型学到的是描述什么画面，尚未学到根据原视频执行修改。图像编辑已有较多“原图—指令—目标图”配对，但把它们逐帧用于视频，相邻帧的编辑会各自独立，杯子的形状可能逐帧跳变。Movie Gen Edit的出发点是：在没有成对视频编辑数据的条件下，借助视频生成模型和若干可构造的代理任务，逐步缩小单帧编辑与真实视频编辑之间的差距\scite{vision-moviegen2024}

主要想法有两点。结构上，把原视频作为额外输入通道接入已经训练好的视频生成主干，使编辑模型继承生成模型的运动先验。监督上，分三个阶段构造训练关系：先用图像编辑配对学单帧修改，再用带运动的合成序列学多帧修改，最后用回译把合成缺陷移到输入侧，让目标回到真实视频。

编辑模型从预训练的文字生成视频模型Movie Gen Video初始化。该模型在时间自编码器（temporal autoencoder, TAE）的潜空间中用流匹配生成视频。TAE在时间、高、宽三个方向各压缩8倍，潜通道为16，潜网格再用核为$1\times2\times2$的三维卷积切成词元。报告给出的最大配置是300亿参数的Transformer（不含文字编码器与TAE），共48层，宽度6144，48个注意力头，前馈层宽度16384，见表\ref{tab:vision-multigen-edit-config}；编辑一节只说明从“预训练的文字生成视频模型”初始化，没有单独写明编辑模型采用哪一个规模。以最大上下文为例，256帧、$768\times768$的视频经TAE压成$32\times96\times96$的潜网格，再按$2\times2$空间块切分，得到$32\times48\times48=73728$个视频词元，约为报告所说的73K。

主干的每一块沿用LLaMa3的Transformer块，使用RMSNorm和SwiGLU前馈层，并作三处修改：注意力改为双向，在自注意力和前馈层之间加入读取文字的交叉注意力，用自适应层归一化接入生成时刻。记第$l$块的输入词元为$\bm Z$、文字表示为$\bm C$、生成时刻嵌入为$\bm e_s$，一块的计算顺序为
\begin{equation}
 \begin{aligned}
  \bm Z'&=\bm Z+\operatorname{SA}\bigl(\operatorname{AdaNorm}(\bm Z;\bm e_s)\bigr),\\
  \bm Z''&=\bm Z'+\operatorname{CA}\bigl(\operatorname{RMSNorm}(\bm Z'),\bm C\bigr),\\
  \bm Z'''&=\bm Z''+\operatorname{SwiGLU}\bigl(\operatorname{AdaNorm}(\bm Z'';\bm e_s)\bigr),
 \end{aligned}
 \label{eq:vision-multigen-moviegen-block}
\end{equation}
其中$\operatorname{SA}$与$\operatorname{CA}$分别是多头自注意力和交叉注意力，$\operatorname{AdaNorm}$表示由$\bm e_s$预测缩放与平移的归一化。报告只说明加入自适应层归一化和交叉注意力的位置，交叉注意力前的归一化形式与调制系数的确切接法这里按常见做法写出，属于示意。位置编码采用可学习的分解形式，分别为时间、高、宽学习嵌入，并在每一层都加入，报告指出这样做减少了时间上的变形。文字由UL2、ByT5和长文本版MetaCLIP三个编码器处理，各自投影到6144维后拼接为$\bm C$，其中按字符编码的ByT5只用于画面中需要出现的文字。

\begin{table*}[tbp]
 \centering
 \small
 \begin{tabularx}{\linewidth}{@{}p{28mm}Xccccc@{}}
  \toprule
  模型 & 输入单元 & 层数 & 宽度 & 头数 & 前馈宽度 & 参数量 \\
  \midrule
  Movie Gen Video & TAE潜网格的$1\times2\times2$块，最多约73K个 & 48 & 6144 & 48 & 16384 & 300亿 \\
  \bottomrule
 \end{tabularx}
 \caption{Movie Gen Edit所依托主干的配置。数值为报告中Movie Gen Video最大规模的配置，参数量只计Transformer本身；编辑模型的规模报告未单独给出。本章各Transformer配置表采用相同的列。}
 \label{tab:vision-multigen-edit-config}
\end{table*}
\FloatBarrier

编辑模型在块嵌入层增加输入通道，把原视频潜网格与含噪目标潜网格沿通道拼接。以教学尺寸为例，64帧、$256\times256$的视频压成$8\times32\times32\times16$的潜网格，两者拼接后为$8\times32\times32\times32$，再切成$8\times16\times16=2048$个词元进入主干。每种编辑任务有一个学习到的任务嵌入：一个线性变换把它映射成四个向量，拼接到文字编码器的隐状态之后；另一个线性变换的输出加到生成时刻嵌入上。所有新增权重初始化为零，其余权重来自预训练的文字生成视频模型，因此训练开始时模型的行为与原生成模型相同。

训练目标仍是流匹配速度回归，只是速度网络多了原视频条件$\bm z^{\mathrm{src}}$和任务嵌入$k$：
\begin{equation}
  L_{\mathrm{edit}}
  =\mathbb E
   \left[
    \left\lVert
     \bm v_{\bm\theta}(\bm z(s),s,c,\bm z^{\mathrm{src}},k)
     -(\bm z_1-\bm z_0)
    \right\rVert_2^2
   \right],
  \label{eq:vision-edit-objective}
\end{equation}
其中$\bm z_1$是目标视频的潜网格，$c$是编辑指令，其余记号同式\eqref{eq:vision-flow-loss}。三个阶段的区别只在于$(\bm z^{\mathrm{src}},c,\bm z_1)$从哪里来。

第一阶段把图像编辑配对视为单帧视频编辑，与文字视频生成交替训练；纯视频生成样本以黑色视频作为空输入条件。论文在约六分之五的批次中采样图像编辑，训练3万步。混合这两类任务，使模型既能利用编辑关系，又不忘记视频运动。单帧训练的时间位置嵌入随机选择，避免所有图像编辑都只出现在序列起点。

单帧训练缺少“原片段连续输入”的经验。第二阶段因此使用两类多帧构造。一类对原图和编辑后图逐步施加相同的仿射变换，形成成对序列。由于两侧使用同样的变换，模型可以学习在相同运动下保留编辑差异，但这类平面移动不能提供真实关节运动。另一类从真实视频定位并分割对象，将相应区域涂成指定颜色，把分割任务转成“把这个对象标成某颜色”的编辑配对，补充真实时间变化。两类构造与文字生成视频按$3/5$、$1/5$、$1/5$的比例采样批次。论文报告这一阶段只训练约1000步，结果变得不那么模糊，但编辑后的元素运动偏少，有时颜色过饱和。

这些目标仍有一部分来自模型生成，可能携带失真。第三阶段使用\term{回译}（backtranslation）：先用前一阶段的模型把真实视频$V$按指令$r$编辑为$\widetilde V$，再由语言模型根据前后描述写出逆向指令$r^{-}$，用$(\widetilde V,r^{-},V)$训练。回译一词借自机器翻译，指把结果翻回原文来构造训练对。这样目标回到真实视频，合成缺陷主要留在输入侧，模型不必把自己的缺陷当作目标学习。例如真实片段里是红花瓶，前向请求为“改成金属杯”，合成输入是带杯子的片段，逆向请求为“把金属杯改回红花瓶”，训练目标仍是原始花瓶视频。这里的$r^{-}$是语义上构造的修改要求，不是保证可逆的数学算子；删除对象可能同时丢失背景细节，逆向指令也可能不准确。报告没有给出这一阶段的训练步数。

推理时，原视频经TAE编码，与指令及任务嵌入一起进入模型，目标潜网格从噪声逐步形成，再解码为编辑片段。检查金属杯的材质之外，还要检查杯体是否在遮挡后恢复为同一形状、玩具狗是否被连带改变。连续编辑时保存每轮输入与结果，才能发现哪一轮开始出现身份或背景漂移。Movie Gen的生成、编辑和配声分别使用相应路线，不能把模型族名称当成一个同时完成全部输出的网络。

\subsubsection{Lumiere视频修补}
\label{sec:vision-multigen-video-inpainting}

文字指令无法精确说明“只改第3至40帧中花瓶占据的像素”。Movie Gen Edit由模型自己推断范围，误改背景时用户没有直接的约束手段。\term{视频修补}（video inpainting）换一种条件：用户给出时空掩码，指出每帧哪些像素允许改变，模型根据其余可见内容和文字补出被遮区域。Lumiere的修补模型沿用前面的时空U-Net，整段时空处理让各帧补出的杯子共享上下文，避免逐帧修补时形状跳变。

结构上只改动输入层。含噪目标视频、保留可见区域的条件视频和二值掩码沿通道拼接：RGB目标3通道、RGB条件3通道、掩码1通道，输入层从3通道扩展到7通道，其余结构与基础模型相同。为便于讨论，本节令$m_{t,i,j}=1$表示允许生成的位置，0表示要保留的位置。训练时从完整视频$\bm X_t$构造条件
$\bm C_t=(\bm1-\bm M_t)\odot\bm X_t$，
掩码沿颜色通道广播；完整视频的含噪版本仍作为待预测对象，损失与式\eqref{eq:vision-lumiere-objective}相同，只是网络多接收了$\bm C_t$与$\bm M_t$。由于掩码可以从任意完整视频随机生成，这种监督不需要人工编辑配对。其他实现可能把掩码语义反过来，接入时须转换，不能把同一个0或1在不同接口中按相同含义解释。

推理时输入原视频、覆盖花瓶的逐帧掩码和替换要求，采样后检查边缘、遮挡及未指定区域。掩码若没有随花瓶移动，模型有时会在旧位置补出杯子，却在新位置留下原花瓶。区域跟踪与掩码传播可利用\ref{chap:vision-video-recognition}章的方法，但传播质量也属于编辑条件的一部分。

若交付明确要求掩码外像素完全相同，可以在制作阶段规定硬合成
\begin{equation}
  \bm X_t^{\mathrm{out}}
  =(\bm1-\bm M_t)\odot\bm X_t+
     \bm M_t\odot\widehat{\bm X}_t\text{。}
  \label{eq:vision-edit-composite}
\end{equation}
该式在掩码为0处代数上保留原像素；这是明确的输出规则，不是模型训练自动提供的保证。硬边界可能暴露接缝，改变光照或物体形状也可能需要更大的编辑范围。若使用软边缘，混合区域就不再逐像素相同，必须相应修改验收范围。

同一接口还能组织动态照片：把一张静态图复制到各个时间位置，只开放局部区域变化。例如让玩具狗摇尾巴而背景固定，掩码应包含尾巴摆动需要扫过的空间。复制帧提供静态背景，没有给出尾巴的唯一真实运动。

两种修改路线的差别见表\ref{tab:vision-video-edit-compare}。Lumiere视频修补用显式掩码换来了明确的修改范围和可自动构造的监督，代价是用户要准备逐帧掩码，而且掩码外的内容只由模型“倾向”保留，严格保持仍需式\eqref{eq:vision-edit-composite}。Movie Gen Edit去掉了掩码，能处理风格替换这类全局修改，代价是监督要分三阶段构造，修改范围依赖模型对指令的理解。

\begin{table*}[tbp]
 \centering
 \small
 \begin{tabularx}{\linewidth}{@{}p{24mm}XX@{}}
  \toprule
  比较项 & Movie Gen Edit & Lumiere视频修补 \\
  \midrule
  修改范围的指定 & 文字指令，由模型推断 & 用户提供逐帧二值掩码 \\
  条件接入 & 原视频潜网格沿通道拼接，任务嵌入接入文字与时刻嵌入 & 条件视频与掩码沿通道拼接，输入层扩展到7通道 \\
  监督来源 & 图像编辑、合成多帧编辑、回译三阶段构造 & 从完整视频随机遮挡自动构造 \\
  适合的修改 & 对象替换、风格变化、全局修改 & 局部区域替换、动态照片 \\
  典型失败 & 误改指令未提及的背景，多轮编辑累积漂移 & 掩码未跟随对象，边缘接缝 \\
  \bottomrule
 \end{tabularx}
 \caption{视频修改两条路线的比较。两者都以原片段为条件；差别在于修改范围由谁指定，以及训练配对从哪里来。}
 \label{tab:vision-video-edit-compare}
\end{table*}
\FloatBarrier

\subsection{视频的延伸与补全}
\label{sec:vision-multigen-extension}
\label{sec:vision-multigen-wan-extension}

延伸与补全的输入是一段或几段已知帧，输出是尚未给定的时间范围。给定历史片段向后续接，只有已发生的状态；给定首尾两侧，中间过程还受到终点限制；给定零散几帧，途中必须经过若干指定状态。难点在于新内容既要与已知帧在边界处衔接，又要在条件不足的方向上给出合理的新动作。创作允许新增动作与事件，而\ref{chap:vision-video-restoration}章的插帧侧重依据已采样观测估计中间时刻；两者可能使用相似掩码，评价证据却不相同。

这三种任务不需要三套模型。Wan的图像条件接口本来就是“已知帧放入条件视频、掩码标出已知位置”的形式，首帧生成只是其中已知帧最少的一种情形。把已知位置换成开头一段，就得到续接；换成首尾两帧，就得到首尾转换；换成零散几帧，就得到其他时间补全。接口的结构与训练见\ref{sec:vision-models-wan}节；报告在预训练中把这些掩码任务混合训练，再对具体任务分别微调。这里只讨论使用。

图\ref{fig:vision-time-masks}使用八个抽样位置展示条件差别。续接给出前三个状态，后面可以朝多种方向发展；首尾条件只固定两端，中间有多种合法运动；零散条件还限制途中必须经过的状态。图中的格子是教学时间位置，不是把八帧直接当作Wan自编码器的合法输入长度。

\begin{figure}[htbp]
  \centering
  \begin{tikzpicture}[x=1cm,y=1cm,font=\figurefont\small,
  cell/.style={draw=BookInk,line width=0.8pt,minimum width=7mm,
    minimum height=7mm,inner sep=0pt},
  known/.style={cell,fill=BookBlue,text=white}]
  \node[anchor=east] at (1.65,2.9) {时间位置};
  \node[anchor=east] at (1.65,2.1) {续接};
  \node[anchor=east] at (1.65,1.1) {首尾转换};
  \node[anchor=east] at (1.65,0.1) {零散条件};
  \foreach \i in {1,...,8} {
    \pgfmathsetmacro{\xx}{1.5+0.95*\i}
    \node at (\xx,2.9) {\i};
    \ifnum\i<4
      \node[known] at (\xx,2.1) {1};
    \else
      \node[cell] at (\xx,2.1) {0};
    \fi
    \ifnum\i=1
      \node[known] at (\xx,1.1) {1};
    \else
      \ifnum\i=8
        \node[known] at (\xx,1.1) {1};
      \else
        \node[cell] at (\xx,1.1) {0};
      \fi
    \fi
    \ifnum\i=2
      \node[known] at (\xx,0.1) {1};
    \else
      \ifnum\i=5
        \node[known] at (\xx,0.1) {1};
      \else
        \ifnum\i=8
          \node[known] at (\xx,0.1) {1};
        \else
          \node[cell] at (\xx,0.1) {0};
        \fi
      \fi
    \fi
  }
  \draw[-{Stealth[length=2mm]},BookInk,line width=0.8pt]
    (2.1,-0.7) -- (9.5,-0.7) node[below left] {时间递增};
\end{tikzpicture}

  \caption{时间条件的三种组织。实心格和数字1表示已提供内容，空心格和数字0表示待生成内容。位置按时间排序，编码与潜空间重排发生在此后；同一掩码形式不保证任意发布权重都支持相应任务。}
  \label{fig:vision-time-masks}
\end{figure}
\FloatBarrier

组织输入时，已知帧按其在目标序列中的时间位置放入条件视频，其余帧填零，再经自编码器编码。掩码保留帧级信息，再按自编码器的时间分组重排到通道维，所以已知帧最好与分组边界对齐：若一个四帧组中只有一帧已知，压缩后的条件潜变量会把已知帧与零帧混在同一组里，模型拿到的是部分信息。续接时通常把上一段末尾若干帧作为新一段的开头。报告构造续接训练数据时，用前1.5秒与后3.5秒画面的SigLIP相似度筛选内容连续的视频；构造首尾转换数据时，则提高首尾差别明显的样本比例。这说明两种任务需要的训练分布不同，使用时应选对应任务微调后的权重。

让玩具狗离开画面后继续前行，模型可以决定离开后的路线；给定它最终位于桌右侧的终帧，就要求生成过程连接两个状态。判断连接质量时，除了比较边界帧外观，还可以记录图像位置的有限差分。若边界前两次横向位移为2、2像素，新片段第一次位移为9像素，即使边界图像接得上，也出现了明显的速度变化。透视和相机运动会影响像素速度，所以此数值只用于固定相机的教学诊断。只给一帧作为续接条件时，模型无法从单帧推断速度，这类速度跳变更常见；给出多帧历史可以提供速度信息。

反复续接时，新片段成为下一轮条件，误差可能进入后续每一段。花瓶位置每次只漂移少量像素，长序列里也可能明显偏离原场景。首尾相接、局部运动连贯与全程状态一致应分别检查。把两端之间的创作画面用于说明一个故事是合理的；把它当成实际未拍到过程的观测证据，则超出了条件所能支持的结论。

\section{视觉条件下的声音生成}
\label{sec:vision-multigen-audio}

本节固定已有视频，待生成的对象是声轨。输入是视频，可选文字说明，有时还有一段已知声音；输出是与视频等长的音频波形。难点有三个：声音内容要符合画面中的声源，即“是什么声音”；起声要落在画面事件上，即“何时发声”；视频可能长达数分钟，而生成网络一次只能处理有限长度。模型可以在训练时同时使用视频、文字与音频，但只要推理输出只改变声音，就仍是视觉条件下的配声任务，与同时创造画面和声音使用不同的可修改范围。

本节按制作步骤展开：先把画面事件和创作选择整理成声音要求，再生成声轨，最后核验并调整。生成步骤中，MMAudio和Movie Gen Audio对上述后两个难点给出不同回答。MMAudio加入专门的高帧率同步特征来解决“何时发声”，处理的是八秒左右的短片段；Movie Gen Audio只用逐帧语义特征，但接收音频上下文，从而能延伸和修补长声轨。

\subsection{声音要求的确定}
\label{sec:vision-multigen-audio-conditions}

画面中一次接触可以对应轻敲、摩擦或破碎，具体声音还取决于材料、力度和录音位置。玩具狗移动也可以配脚步或机械声。视频往往没有唯一决定这些声学属性，因此需要把可见事件与创作选择分开记录。

可以用事件记录$(\tau_i^{\mathrm{start}},\tau_i^{\mathrm{end}},e_i,r_i)$描述第$i$个声音要求，其中时间以秒计，$e_i$表示对应事件，$r_i$说明声源、内容及保持要求。某次敲击记录来自视频中的接触；画外鸟鸣记录来自创作说明。两者都能成为有效条件，但后者不是模型从像素直接测得的事实。

同一四秒教学片段中，假定两次接触发生于1.0秒和2.5秒，可规定每次接触配一个短促轻响，其他时间保留低音量环境声。这份记录不等于模型一定支持逐事件数值控制；如果入口只有文字和视频，它只是验收依据。只有模型实际接收时间标记、局部掩码或已知音频上下文时，这些字段才直接进入生成计算。

\subsection{声轨的生成}
\label{sec:vision-multigen-audio-alignment}

语义特征帮助决定“什么声源”，时间特征帮助决定“何时发声”。一段视频的平均特征可以保留花瓶和桌面，却难以说明两次敲击之间的停顿；每秒只取几帧的语义特征，也可能恰好错过只持续几十毫秒的接触瞬间。生成器必须同时利用场景内容与局部变化，才能把声轨安排到画面事件上。

两个模型都在音频潜空间中用流匹配生成，也都把视觉特征按时间对齐到音频潜序列。它们的差别在于时间信息的来源和长度的处理。MMAudio在每秒8帧的语义特征之外，另加每秒24帧、专门为音画同步训练的特征；Movie Gen Audio的视觉特征只有逐帧语义嵌入，但它把已知音频作为上下文条件，用同一个模型完成整段生成、延伸和修补。

\subsubsection{MMAudio}
\label{sec:vision-multigen-mmaudio}

视频配声的配对数据有限，常用的VGGSound约有18万段十秒视频，MMAudio训练时取每段前8秒；只用这些数据训练，模型对声音种类和文字描述的覆盖不足。一种做法是先训练文字生成音频模型，再加入视觉条件微调，但这样视觉和文字在不同阶段学习，二者之间的交互较弱。MMAudio的主要想法是从头联合训练：同一个网络同时接收音视频配对和音频文字配对，缺失的模态用学习到的空条件替代，这样大规模音频文字数据也能帮助视频配声\scite{vision-mmaudio2025}。为了解决时间对应，它再加入帧级同步特征，以逐位置调制的方式影响音频生成。

音频先转为梅尔频谱，再由预训练的自编码器压成潜序列。梅尔频谱的计算已在\ref{sec:vision-models-encoding}节说明：把波形切成重叠短窗，对每窗求频谱，再按接近听觉尺度的梅尔刻度分组汇总能量，得到“时间窗$\times$频率组”的二维数组；这里只需记住它的一维是时间，所以压缩后的潜序列仍按时间排列。默认16 kHz版本的潜序列为每秒31.25个位置、每个位置20维；44.1 kHz版本为每秒约43.07个位置、40维。视觉帧和文字经CLIP提取特征，视觉特征为每秒8帧、1024维，文字为77个词元、1024维。生成主干先有$N_1$个多模态Transformer块，再接$N_2$个只更新音频表示的单模态块，最后预测音频潜变量的速度。

多模态块大体沿用SD3的MM-DiT设计：文字、视觉、音频三条序列各自保留一条计算支路，只在注意力处汇合。设三条序列在某一块中的查询、键、值分别为$(\bm Q_{\mathrm t},\bm K_{\mathrm t},\bm V_{\mathrm t})$、$(\bm Q_{\mathrm v},\bm K_{\mathrm v},\bm V_{\mathrm v})$、$(\bm Q_{\mathrm a},\bm K_{\mathrm a},\bm V_{\mathrm a})$，联合注意力把它们沿序列方向拼接后做一次注意力，
\begin{equation}
 [\bm O_{\mathrm t};\bm O_{\mathrm v};\bm O_{\mathrm a}]
 =\operatorname{softmax}\!\left(
   \frac{[\bm Q_{\mathrm t};\bm Q_{\mathrm v};\bm Q_{\mathrm a}]
         [\bm K_{\mathrm t};\bm K_{\mathrm v};\bm K_{\mathrm a}]^\top}{\sqrt{d_{\mathrm h}}}
  \right)[\bm V_{\mathrm t};\bm V_{\mathrm v};\bm V_{\mathrm a}],
 \label{eq:vision-multigen-joint-attention}
\end{equation}
再把输出按原长度切回三条序列，其中$[\cdot;\cdot]$表示沿序列方向拼接，$d_{\mathrm h}$为每个注意力头的维数。这样每个音频位置可以同时读取文字、视觉和其他音频位置。块内其余部分是：由全局条件经自适应层归一化提供的逐通道缩放与偏移；视觉和音频支路用核大小为3的一维卷积代替前馈层中的线性层（论文称为ConvMLP），使每个位置还能看到相邻时间位置；文字支路仍用普通前馈层。单模态块去掉文字与视觉支路，联合注意力退化为音频自注意力。论文给出的四种规模见表\ref{tab:vision-multigen-audio-config}，默认小模型取$N_1=4$、$N_2=8$、隐藏宽度448，约1.57亿参数；最大的L版本把块数加到7和14、宽度加到896，约10.3亿参数。论文没有给出注意力头数和前馈层宽度。

三种序列不具有相同的时间分辨率。音频潜序列为每秒31.25个位置，视觉语义特征为每秒8个位置，把二者相同的序号直接当作同一时刻会发生错位。模型在联合注意力之前对视觉和音频支路的查询与键施加旋转位置编码，并把视觉支路的频率乘以$31.25/8$：第$k$个视觉位置的有效位置变为$31.25k/8$，例如第8个视觉位置（第1秒）对应有效位置31.25，正好与第1秒的音频位置相同。文字没有时间含义，不加旋转位置编码。

更细的同步信息来自在音画偏移检测上预训练的Synchformer视觉编码器。它读取短时视频，提供与事件发生时间相关的特征，有效特征率为每秒24个位置、768维；附录的提取过程从每秒25帧输入形成重叠短片段，不能把24直接当作所有输入长度下不变的原始取样率。同步特征经过局部卷积处理，并用最近邻插值对齐到音频潜位置，再加上由文字、视觉及生成时刻得到的全局条件。令对齐后的第$j$个条件为$\bm c_j$，该位置音频隐状态为$\bm h_j\in\mathbb R^d$，调制写成
\begin{equation}
  \widetilde{\bm h}_j
  =\bm\gamma(\bm c_j)\odot
     \operatorname{LayerNorm}(\bm h_j)
    +\bm\beta(\bm c_j)\text{。}
  \label{eq:vision-audio-sync-modulation}
\end{equation}
两个学习映射给出逐通道缩放与偏移。全局条件会在所有$j$处重复，同步条件则随$j$变化，因此第1秒附近可以得到与第2.5秒不同的调制。这里没有把同步编码器的偏移分类结果直接当作波形，它提供的是生成条件。

在默认设置中，八秒音频包含128000个波形采样点，压缩后有250个潜位置；同一段视频的语义特征有64个位置，同步特征有192个位置。对齐不是要求这三个长度相等，而是让它们按共同时间关系进入交互。插值不会新增视频中未观察到的敲击事件，特征本身漏掉动作时，调制也不能凭空恢复证据。

训练数据包括VGGSound的音视频配对，以及AudioCaps、Clotho和WavCaps等音频文字数据，后者合计约95万对。训练从随机初始化开始，把真实音频潜变量与噪声连接，回归条件速度，复用式\eqref{eq:vision-flow-loss}的目标。视觉条件和文字条件各以10\%的概率替换为空条件；这样同一网络既学会了带条件的速度，也学会了不带某个条件的速度，推理时可以做无分类器引导。音频文字数据本来没有视频，视觉与同步入口始终使用空条件。

推理固定视频和文字，从音频噪声出发用欧拉法积分25步，无分类器引导强度取4.5，得到干净潜表示，再由自编码器恢复频谱、由\term{声码器}（vocoder）恢复波形。声码器是把频谱转换成可播放波形采样的网络，它不负责判断画面中的碰撞是否发生。论文报告小模型生成八秒声音约需1.23秒。固定花瓶落桌片段，分别给出“轻轻敲击”和“陶瓷破碎”的文字，可以区分内容选择与时间对应。若画面没有破碎而声轨出现碎裂声，作品可能满足了一个夸张音效要求，却不符合写实碰撞要求。MMAudio的音效和环境声结果也不能证明长段对白可懂、说话人稳定或歌词准确，这些输出需要自己的训练范围与评价。

\subsubsection{Movie Gen Audio}
\label{sec:vision-multigen-movie-audio}

MMAudio一次处理八秒左右，长片段若逐段独立生成，相邻段的背景声、响度和混响可能突变；它也没有接收已知声音的入口，不能只替换声轨中的一小段。Movie Gen Audio的主要想法是把已知音频作为上下文条件：训练时随机遮去目标音频的一部分，让模型根据视频、文字及其余已知声音预测被遮部分\scite{vision-moviegen2024}。没有上下文对应整段生成，已知一侧对应延伸，两侧已知对应局部修补，长声轨则通过反复延伸或多窗口协调形成。

模型是约130亿参数的扩散Transformer，用流匹配训练，输出48 kHz音频。波形先由单独训练的DAC-VAE压成每秒25帧、每帧128维的连续潜序列。上下文音频同样编码，保留内容的位置作为条件，其余位置填零，再与含噪目标潜变量逐位置沿通道拼接。视频每帧经长提示版MetaCLIP得到1024维嵌入，每个音频潜位置取时间上最近的视频帧，投影后逐位置加到音频表示上，而不是沿时间拼接成更长的序列。文字经T5-base编码为最多512个词元的768维特征，在每个Transformer块中通过交叉注意力接入。以十秒片段为例，目标潜序列为$250\times128$，与上下文拼接后每个位置有256维输入，另有250个对齐的视觉嵌入。

主干共36层，宽度4608，前馈层宽度18432，即四倍于宽度；130亿参数不含MetaCLIP、T5和DAC-VAE，注意力头数论文未给出。每一块依次是自注意力、读取文字的交叉注意力和前馈层，生成时刻的接入方式与\ref{sec:vision-models-wan}节回顾的Wan相同：一个所有层共享的小网络读取生成时刻嵌入，预测四个缩放和两个偏置共六个调制参数，每层再加上自己的一组偏置。把式\eqref{eq:vision-multigen-moviegen-block}中的自适应归一化理解为这六个参数对自注意力与前馈层输入输出的缩放平移，就得到这里一块的计算顺序。视频嵌入经带门控的线性投影后，与音频表示逐位置相加。表\ref{tab:vision-multigen-audio-config}把它与MMAudio的四种规模放在同一组列下。

\begin{table*}[tbp]
 \centering
 \small
 \begin{tabularx}{\linewidth}{@{}p{28mm}Xccccc@{}}
  \toprule
  模型 & 输入单元 & 层数 & 宽度 & 头数 & 前馈宽度 & 参数量 \\
  \midrule
  MMAudio S-16kHz & 每秒31.25个音频潜位置，20维 & 4+8 & 448 & 未给出 & 未给出 & 1.57亿 \\
  MMAudio S-44.1kHz & 每秒约43.07个音频潜位置，40维 & 4+8 & 448 & 未给出 & 未给出 & 1.57亿 \\
  MMAudio M-44.1kHz & 同上 & 4+8 & 896 & 未给出 & 未给出 & 6.21亿 \\
  MMAudio L-44.1kHz & 同上 & 7+14 & 896 & 未给出 & 未给出 & 10.3亿 \\
  Movie Gen Audio & 每秒25个DAC-VAE潜位置，128维，与上下文拼接为256维 & 36 & 4608 & 未给出 & 18432 & 130亿 \\
  \bottomrule
 \end{tabularx}
 \caption{两种视觉条件声音生成模型的配置。MMAudio的层数写成“多模态块数+单模态块数”；“未给出”表示论文没有报告该项。Movie Gen Audio的参数量不含文字、视频编码器与DAC-VAE。}
 \label{tab:vision-multigen-audio-config}
\end{table*}
\FloatBarrier

训练先在约百万小时量级的音频上预训练，单次长度不超过30秒，再用带视频的高质量数据微调。目标仍是式\eqref{eq:vision-flow-loss}形式的速度回归，只是速度网络多了视频、文字和音频上下文条件。论文明确说明模型不覆盖语音和带人声的音乐。推理用中点法求解64步，无分类器引导权重为3，并为每条请求生成8个候选再重排挑选。一次最多生成约30秒，即750个潜位置，更长的声轨用两种延伸方式。一种逐段自回归：把上一段末尾的音频作为下一段的已知上下文。若两个片段都包含同一次碰撞，就不能在最终结果里各保留一次，必须按共同时间位置合并重叠区域。另一种是多窗口共同更新：每次生成步把全局音频潜序列分成重叠窗口，各自预测，再按三角形窗口权重归一化合并，之后进入下一步。论文报告后者略好，作为默认方式。

多窗口共同更新就是算法\ref{alg:vision-multigen-window-consensus}，只是把\ref{sec:vision-multigen-lumiere}节的等权平均换成三角形加权平均。用中间高、两端低的三角形权重，窗口边缘的预测在合并中占比较小，可以减少从单窗口突然切到等权混合造成的突变；若某个位置的权重总和为零，则无法归一化，窗口布置必须保证覆盖。论文用窗口长度$n_{\mathrm{win}}$、相邻窗口终点之差$n_{\mathrm{hop}}$和重叠长度$n_{\mathrm{ctx}}=n_{\mathrm{win}}-n_{\mathrm{hop}}$描述窗口布置，以帧为单位，默认取$40$、$30$和$10$；在音效与音乐延伸实验中，它又以秒为单位取$n_{\mathrm{hop}}=n_{\mathrm{ctx}}=5.5$秒。

教学上，先为固定视频形成环境声，再只替换第二次敲击。条件声轨在第二次敲击附近留空，其余部分作为上下文；生成后检查新敲击是否出现、周围环境是否连续。即使潜空间有上下文条件，解码器仍可能影响邻近波形，所以要求严格保持的区间要在最终波形上核验。

两个模型的比较见表\ref{tab:vision-audio-gen-compare}。Movie Gen Audio相对MMAudio改变了两处设计：去掉专门的同步特征，改用每秒25帧的逐帧语义嵌入；增加音频上下文入口。前者使时间对应依赖模型从语义帧的变化中学习，后者换来了延伸、修补和长声轨能力。代价是模型规模大一到两个数量级，而且对快速连续的短促事件，没有专门同步特征时的对应精度需要单独核验。

\begin{table*}[tbp]
 \centering
 \small
 \begin{tabularx}{\linewidth}{@{}p{24mm}XX@{}}
  \toprule
  比较项 & MMAudio & Movie Gen Audio \\
  \midrule
  规模与输出 & 1.57亿至10.3亿参数，16 kHz或44.1 kHz & 约130亿参数，48 kHz \\
  语义条件 & CLIP视觉特征每秒8帧，联合注意力接入 & MetaCLIP逐帧嵌入，按最近帧对齐后相加 \\
  时间对应 & Synchformer同步特征每秒24帧，逐位置调制 & 无专门同步特征，依赖逐帧嵌入 \\
  长度与上下文 & 约八秒片段，无音频上下文入口 & 单次约30秒；音频上下文支持延伸与修补 \\
  训练 & 音视频与音频文字数据从头联合训练 & 大规模音频预训练，再带视频微调 \\
  不覆盖 & 长段对白、歌词等未作为目标 & 语音与带人声的音乐 \\
  \bottomrule
 \end{tabularx}
 \caption{两种视觉条件声音生成模型的比较。数值来自各自论文的说明，不是同一测试集上的质量比较。}
 \label{tab:vision-audio-gen-compare}
\end{table*}
\FloatBarrier

\subsection{声轨的核验}
\label{sec:vision-multigen-audio-check}

声音可以“听起来像敲击”，却早于画面半秒出现。对短促事件，可以先匹配声源与次数，再比较已匹配事件的起始时刻。设视频接触时刻为$\tau_i^{\mathrm v}$，音频起声时刻为$\tau_i^{\mathrm a}$，定义
\begin{equation}
  \delta_i=\tau_i^{\mathrm a}-\tau_i^{\mathrm v},
  \qquad
  e_{\mathrm{sync}}=\frac1m\sum_{i=1}^{m}|\delta_i|\text{。}
  \label{eq:vision-event-sync-error}
\end{equation}
$\delta_i>0$表示声音较晚。这里的$m$是已正确配对事件数，漏声、多声与声源错误另行计数，不能只对少量成功配对求平均而隐藏其余失败。

计算偏移之前要先回答“哪次起声对应哪次接触”。这正是\ref{sec:vision-alignment-alignment}节讨论的跨模态对齐问题，只是对齐的粒度是事件，而不是逐帧或整段。事件多、间隔不均匀时，可以借用\ref{sec:vision-alignment-dtw}节的动态时间规整：把接触时刻序列与起声时刻序列看成两条序列，逐对代价取时间差的绝对值，在保持先后顺序的前提下求累积代价最小的对应路径。路径的单调性保证了第一次接触不会与第二次起声配对、第二次接触再与第一次起声配对。但动态时间规整要求每个元素都有对应，它会把多出的一声强行并到某次接触上；所以在配对之后，应当把时间差超过事先规定的匹配窗口的对应拆开，计为多声或漏声；匹配窗口回答“算不算同一事件”，比下文判断同步是否合格的容差宽。教学上，若接触为1.0秒和2.5秒，起声为1.12秒、1.30秒和2.48秒，规整路径会把前两次起声都对应到第一次接触。取0.25秒的匹配窗口，1.30秒那一声与接触相差0.30秒，应计为多出的声音，而不是第一次接触的另一个偏移；1.12秒那一声保留为配对，再进入下面的同步判断。

沿用1.0秒和2.5秒的接触，若起声时刻为1.12秒和2.48秒，偏移为120毫秒和$-20$毫秒，平均绝对偏移70毫秒。若教学验收预先规定每次不得超过80毫秒，则只有第二次通过。这是便于计算的任务阈值，不是普适的人类听觉阈值。图\ref{fig:vision-event-sync}把两种方向的偏移放到共同时间轴上。

\begin{figure}[htbp]
  \centering
  \begin{tikzpicture}[x=1cm,y=1cm,font=\figurefont\small,
  axis/.style={BookInk,line width=0.8pt},
  contact/.style={BookBlue,line width=1.2pt},
  sound/.style={BookAmber,line width=1.2pt}]
  \node[anchor=east] at (1.6,3.4) {接触};
  \node[anchor=east] at (1.6,2.6) {起声};
  \draw[axis] (2,3.4) -- (10,3.4);
  \draw[axis] (2,2.6) -- (10,2.6);
  \foreach \t in {0,...,4} {
    \pgfmathsetmacro{\xx}{2+2*\t}
    \draw[axis] (\xx,3.35) -- (\xx,3.45);
    \draw[axis] (\xx,2.55) -- (\xx,2.65);
    \node[below] at (\xx,2.45) {\t};
  }
  \node[anchor=west] at (7.45,1.78) {媒体时间（秒）};
  \foreach \xx in {4,7} {
    \draw[contact] (\xx,3.2) -- (\xx,3.65);
    \fill[BookBlue] (\xx,3.4) circle (1.5pt);
  }
  \foreach \xx in {4.24,6.96} {
    \draw[sound] (\xx,2.6) -- ++(-0.1,-0.16) -- ++(0.2,0) -- cycle;
  }
  \node at (4,3.95) {1.00};
  \node at (7,3.95) {2.50};
  \node[anchor=west] at (0.3,1.3) {第一次：声音较晚};
  \node[anchor=west] at (5.75,1.3) {第二次：声音较早};
  \draw[axis,-{Stealth[length=2mm]}] (0.6,0.1) -- (5,0.1);
  \draw[axis,-{Stealth[length=2mm]}] (6,0.1) -- (10.4,0.1);
  \draw[contact] (1.7,-0.08) -- (1.7,0.35);
  \draw[sound] (3.9,-0.08) -- (3.9,0.35);
  \draw[sound] (7.4,-0.08) -- (7.4,0.35);
  \draw[contact] (8.35,-0.08) -- (8.35,0.35);
  \draw[BookInk,-{Stealth[length=2mm]},line width=0.8pt]
    (1.7,0.58) -- (3.9,0.58);
  \draw[BookInk,-{Stealth[length=2mm]},line width=0.8pt]
    (8.35,0.58) -- (7.4,0.58);
  \node[above] at (2.8,0.62) {$+120$ ms};
  \node[above] at (7.88,0.62) {$-20$ ms};
  \node[below] at (1.7,-0.12) {1.00};
  \node[below] at (3.9,-0.12) {1.12};
  \node[below] at (7.25,-0.12) {2.48};
  \node[below] at (8.5,-0.12) {2.50};
  \node[anchor=west] at (0.3,-0.85) {局部放大分别使用自己的时间尺度，刻度单位为秒。};
\end{tikzpicture}

  \caption{两次碰撞的教学时间对齐。上排为可见接触，下排为音频起声，局部放大给出偏移方向与大小；不是模型实测波形。平均值可能掩盖第一次事件超出既定容差，漏声和多声还须单独记录。}
  \label{fig:vision-event-sync}
\end{figure}
\FloatBarrier

逐事件标注成本较高时，可以使用自动估计。MMAudio论文用DeSync指标评价同步：由Synchformer预测整段音画的时间偏移，取其绝对值，单位为秒。这种估计不需要事件标注，但它给出的是整段的单一偏移，不能指出哪一次碰撞出了问题；而且评价用的Synchformer与MMAudio生成时用的同步编码器相同，模型可能恰好迎合了该编码器的偏好。因此自动指标适合在模型之间做大批量比较，最终交付仍应对关键事件做式\eqref{eq:vision-event-sync-error}的逐事件检查。

持续环境声没有唯一冲击起点，应看场景与前后连贯性；对白还需要转写内容、说话人和口型关系。近景同步测量也要说明是否扣除音频设备延迟，远距离声源则可能具有合理的传播延迟。声画不同时出现并不总是错误，判断取决于所描述的物理或创作条件。

调整要针对失败来源。声源不合适可改文字，局部时间错误可用上下文修补重做相应区间，长背景突变可改用上下文延伸或多窗口协调。每次保存改动范围和候选，重新检查与邻段的连接。修改、延伸与修补是不同条件下的选择，不是所有配声都必须经历的固定步骤。

\section{音视频的协同生成}
\label{sec:vision-multigen-joint}

当画面与声音都尚未固定，输入只有文字（有时加一幅首帧），输出是一段带声轨的视频。难点在于两种媒体要描述同一组事件：玩具狗碰到花瓶的那一帧，正是碰撞声开始的时刻；人物开口说话时，嘴型要对应发出的音节。前一节的配声模型只能让声音去适应已经固定的画面，如果画面里的碰撞本身不清楚，声音无法补救。

两条路线对“画面和声音在什么时候互相约束”给出不同回答。分阶段条件传递先生成一种媒体，再以它为条件生成另一种，约束是单向的。联合采样把两种媒体放进同一个生成过程，每一步都让双方读取对方的当前状态，约束是双向的。后一条路线由Ovi和LTX-2实现，二者结构相近，差别在于训练组织和推理时的引导方式。

\subsection{基于分阶段条件传递的生成}
\label{sec:vision-multigen-staged}

先制作玩具狗碰倒花瓶的视频，再按成片配声，可以直接复用前两节的模型。视频阶段交付的不只是帧文件，还应包括实际时长、帧率、事件时间、可见声源以及声音可自由设计的部分。若碰撞从原计划的1.0秒变成1.4秒，配声条件应采用成片中的时间。

制作过程中可把理想化的分布关系写成
\begin{equation}
  p(V,A\mid c)=p(V\mid c)\,p(A\mid V,c),
  \label{eq:vision-staged-factorization}
\end{equation}
其中$V$为视频，$A$为音频，$c$为创作条件。这是概率链式分解，本身没有损失一般性；实际误差来自两个模型的近似、条件传递以及第二阶段受已选视频约束。不能仅从分阶段形式推出它在数学上无法表示联合分布。问题出在第一阶段：$p(V\mid c)$的模型在生成时并不知道之后要配什么声音，它可能生成一个碰撞不清楚的画面，而第二阶段没有办法修改它。

在固定成片上配声时，声音模型只能适应已有碰撞。若画面没有真正接触花瓶，补一声碰撞不会让动作成立。此时应回到视频阶段修改，再更新配声，或者选择允许共同调整的模型。分阶段路线的好处是每个阶段都可以选用最合适的专用模型，并能单独替换和核验；不同模型间共享时间单位和事件记录，可以组成一致的制作流程，但不等于它们共享参数或共同训练。说话场景是这条路线最难处理的情形：先生成的人物嘴部运动没有对应的语音内容，后配的语音很难与嘴型吻合。

\subsection{基于联合采样的生成}
\label{sec:vision-multigen-joint-mechanism}

联合采样把画面和声音都放在待采样状态中。文字应同时说明可见动作与声音要求，训练配对要来自同一时间范围，否则分支之间学到的联系可能是错误的。两个代表模型都采用“每种媒体一个分支、分支之间双向交叉注意力”的结构。Ovi的两个分支结构完全相同，音频分支先单独预训练，再与视频分支融合；LTX-2的两个分支宽度可以不同，结构和训练已在\ref{sec:vision-models-ltx2}节说明，这里只讲它在推理时怎样分别调节文字和另一模态的影响。

\subsubsection{Ovi}
\label{sec:vision-multigen-ovi}

分阶段路线无法让先生成的画面考虑声音。若直接训练一个大网络同时输出两种潜变量，视频和音频的潜空间、序列长度都不同，从头训练还要同时学会两种媒体的生成，成本很高。Ovi的主要想法是让两个分支采用完全相同的结构：视频分支取自已经训练好的视频生成模型，音频分支用同样的结构单独学会生成声音，然后在每一层用双向交叉注意力把两者连接起来，只训练连接和注意力部分\scite{vision-ovi2025}。结构对称使两个分支每一层的隐藏宽度相同，交叉注意力可以直接逐层配对。

视频分支由Wan2.2的50亿参数模型初始化，音频分支结构相同、从头训练，合计约110亿参数。每个分支有30个块，隐藏宽度3072，24个注意力头、每头128维，前馈层宽度14336，配置见表\ref{tab:vision-multigen-ovi-config}。每块包含自注意力、文字交叉注意力、音视频交叉注意力和前馈层。论文列出了每块的这几类子层，但没有写明它们的先后顺序和生成时刻的接入方式。为了说明数据怎样流动，下面按“先自注意力、再两种交叉注意力、最后前馈层”的示意顺序写出视频分支第$l$块，记视频词元为$\bm Z^{\mathrm v}$、同一块中音频分支的隐藏表示为$\bm Z^{\mathrm a}$、文字表示为$\bm C$：
\begin{equation}
 \begin{aligned}
  \bm Z_1&=\bm Z^{\mathrm v}+\operatorname{SA}(\operatorname{Norm}(\bm Z^{\mathrm v})),\\
  \bm Z_2&=\bm Z_1+\operatorname{CA}(\operatorname{Norm}(\bm Z_1),\bm C),\\
  \bm Z_3&=\bm Z_2+\operatorname{CA}(\operatorname{Norm}(\bm Z_2),\bm Z^{\mathrm a}),\\
  \bm Z_4&=\bm Z_3+\operatorname{FFN}(\operatorname{Norm}(\bm Z_3)),
 \end{aligned}
 \label{eq:vision-multigen-ovi-block}
\end{equation}
音频分支对称地以视频表示作为音视频交叉注意力的键和值。式中的顺序与归一化位置是示意，不是论文给出的配置。两个分支宽度相同，第三行的交叉注意力不需要额外的维度投影就能逐层配对；30个块各有一个音视频交叉注意力，所以两种媒体在整个网络深度上反复交换信息。

\begin{table*}[tbp]
 \centering
 \small
 \begin{tabularx}{\linewidth}{@{}p{28mm}Xccccc@{}}
  \toprule
  模型 & 输入单元 & 层数 & 宽度 & 头数 & 前馈宽度 & 参数量 \\
  \midrule
  Ovi视频分支 & Wan2.2视频潜网格的时空块，121帧对应31个潜时间位置 & 30 & 3072 & 24 & 14336 & 约50亿 \\
  Ovi音频分支 & MMAudio 16 kHz自编码器潜序列，五秒157个位置 & 30 & 3072 & 24 & 14336 & 约50亿 \\
  \bottomrule
 \end{tabularx}
 \caption{Ovi两个分支的配置。论文给出全模型约110亿参数、每个分支约50亿参数，两个分支之和与全模型规模之间的差额论文未逐项说明；融合阶段约57亿参数参与更新。}
 \label{tab:vision-multigen-ovi-config}
\end{table*}
\FloatBarrier

输入的视频是121帧、每秒24帧、约$720\times720$像素的片段，约5秒；视频自编码器把它压成31个潜时间位置。音频用MMAudio的16 kHz自编码器，五秒音频按每512个采样一个位置，得到157个潜位置；生成后先解码为频谱，再由BigVGAN声码器恢复波形。文字由一个冻结的T5编码器处理，两个分支共享同一条文字条件。提示中用\code{<S>}与\code{<E>}包住要说出的台词，用\code{<AUDCAP>}与\code{<ENDAUDCAP>}包住对声音的描述，放在提示末尾。

交叉注意力读取另一模态的隐藏特征，需要知道哪些位置对应同一时刻。视频有31个时间位置，音频有157个，二者覆盖同样五秒。若直接用序号作旋转位置编码，第31个音频位置会被当作与第31个视频位置同时，而前者实际只在约1秒处。Ovi把音频的旋转位置编码频率乘以$31/157\approx0.197$：第$j$个音频位置的有效位置变为$0.197j$，第157个位置约为31，与最后一个视频位置对齐。这一缩放只建立时间尺度的对应，正确的敲击同步仍要从配对数据中学习。

训练分三个阶段。第一阶段只训练音频分支：用数十万小时、以语音为主的音频，长度最多12秒，学习文字到声音的生成；论文报告训练5万步，批量2880，学习率$10^{-4}$。第二阶段把音频长度固定为约5秒，并加入音效数据微调，使音频分支与视频片段长度一致。第三阶段融合：新建的音视频交叉注意力从随机初始化开始，两个分支的前馈网络全部冻结，只训练自注意力和交叉注意力，110亿参数中约57亿参与更新，训练4万步，批量768，学习率$5\times10^{-5}$。配对音视频提供两种干净目标$\bm z_1^{\mathrm v}$和$\bm z_1^{\mathrm a}$，各自抽取独立噪声，但使用同一个生成时刻$s$。损失是两个流匹配目标的加权和：
\begin{equation}
  L_{\mathrm{joint}}
  =w_{\mathrm v}\left\lVert\bm v^{\mathrm v}-(\bm z_1^{\mathrm v}-\bm z_0^{\mathrm v})\right\rVert_2^2
   +w_{\mathrm a}\left\lVert\bm v^{\mathrm a}-(\bm z_1^{\mathrm a}-\bm z_0^{\mathrm a})\right\rVert_2^2,
  \label{eq:vision-ovi-objective}
\end{equation}
其中$w_{\mathrm v},w_{\mathrm a}$由论文设定。共享$s$的原因是推理时两个分支要沿同一日程前进；若训练时两边噪声水平不同，交叉注意力学到的是“干净程度不同的两种媒体”之间的关系，与推理时的情况不符。

推理时，把联合状态视为一个较长的向量，其演化可写为
\begin{equation}
 \frac{\mathrm d}{\mathrm ds}
 \begin{pmatrix}\bm z^{\mathrm v}\\ \bm z^{\mathrm a}\end{pmatrix}
 =
 \begin{pmatrix}
 \bm v^{\mathrm v}(\bm z^{\mathrm v},\bm z^{\mathrm a},s,c)\\
 \bm v^{\mathrm a}(\bm z^{\mathrm a},\bm z^{\mathrm v},s,c)
 \end{pmatrix}\text{。}
 \label{eq:vision-joint-flow}
\end{equation}
每边速度都依赖两边当前状态。采用欧拉法说明时，两个下一步都应由同一个当前状态计算，再共同更新；若先把视频推进一步，然后用新视频与旧音频计算音频速度，就已换成另一种分步更新。论文用UniPC求解器共同积分两个分支，并报告它比欧拉法稳定；无分类器引导在每一步增加一次去条件前向。

教学上取单维潜状态$z^{\mathrm v}=0.2$、$z^{\mathrm a}=-0.1$，假设当前预测速度分别为0.6和0.4，步长0.25。一次共同欧拉更新得到0.35和0。这个计算只说明数值更新的同步，没有证明媒体里的敲击已同步。生成时刻$s$走到1以后，还需分别解码视频和音频，按媒体时间组合播放。

请求中可以同时规定“玩具狗接触花瓶”“短促金属碰撞声”和“安静房间环境”。结果应检查声源、动作及起止是否属于同一事件。说话场景还要检查台词、音素和嘴部运动，不能用敲击同步代替口型验证。Ovi原版聚焦约五秒片段；将多段分别成功的片段拼成一分钟故事，仍需另外解决跨镜头身份、场景与叙事状态。

\subsubsection{LTX-2}
\label{sec:vision-multigen-ltx}

LTX-2的视频与音频分支、双向交叉注意力及联合解码已在\ref{sec:vision-models-ltx2}节说明，两个分支的层数与宽度等配置以该节为准：视频和音频各由自己的自编码器压成潜表示，两个分支通过双向交叉注意力交换信息，跨模态连接主要使用时间位置，训练用流匹配监督两路速度，文字条件来自冻结的Gemma3隐表示\scite{vision-ltx22026}。与Ovi相同，它在每个生成步共同更新两路状态。这里讨论使用时的一项差别：推理时可以分别增强文字条件和另一模态带来的影响，使用差异发生在每一步的预测组合中，不是生成完成后才为声音与画面打分。

固定当前模态的潜状态及生成时刻，记$\bm f(c,m)$为同时给定文字$c$和另一模态信息$m$时的预测，$\bm f(\varnothing,m)$为去掉文字的预测，$\bm f(c,\varnothing)$为去掉另一模态信息的预测。报告的引导组合可写为
\begin{equation}
 \widehat{\bm f}
 =\bm f(c,m)
  +\lambda_c[\bm f(c,m)-\bm f(\varnothing,m)]
  +\lambda_m[\bm f(c,m)-\bm f(c,\varnothing)]\text{。}
 \label{eq:vision-modality-guidance}
\end{equation}
两项差分别刻画加入文字与加入另一模态后当前预测如何变化。这里把完整条件预测作为基准，故$\lambda_c=\lambda_m=0$仍保留两种条件，不能误读为无条件生成。去条件分支也应使用方法规定的缺省入口，不是任意填一幅黑图就等价。每一步因此需要三次前向，计算量约为不加引导时的三倍。

假设某个潜分量三次预测依次为0.6、0.2、0.5，取$\lambda_c=2$、$\lambda_m=1$，组合值为$0.6+2(0.6-0.2)+(0.6-0.5)=1.5$。它是潜空间预测，可以超出三次预测的范围，不是必须在0至1之间的像素。提高引导会改变采样轨迹，过大时画面可能过饱和或动作僵硬，是否改善当前请求需比较完整结果。

调节时，可以在同一“玩具狗靠近花瓶并发出机械声”请求下固定初始噪声、长度与采样步数，只改变$\lambda_m$，查看声音是否更贴近动作、画面是否出现额外变化。随后再在多组噪声下重复，避免单个候选主导判断。若文字与首帧冲突，增强两个条件也不自动解决冲突，需要先明确优先保留哪一个。高分辨率制作可以先生成基础潜变量，再做空间潜上采样，并对重叠时空区块细化后合并。两路潜变量最终分别恢复为视频帧与音频频谱，频谱经声码器形成波形。容器层仍要写入正确时长、采样率与时间戳；如果导出时错误改变帧率，模型内部的时间联系也会被破坏。分块降低资源压力，不提供无限长叙事的一致性保证。

三种协同方式的比较见表\ref{tab:vision-joint-compare}。Ovi相对分阶段路线把单向约束改成每步的双向读取，换来了画面能考虑声音、说话场景能联合生成嘴型与语音，代价是需要大规模配对音视频数据和双倍的主干计算。LTX-2保留了联合采样，又在推理时把两种条件的强度拆开调节，代价是每步三次前向。

\begin{table*}[tbp]
 \centering
 \small
 \begin{tabularx}{\linewidth}{@{}p{22mm}XXX@{}}
  \toprule
  比较项 & 分阶段条件传递 & Ovi & LTX-2 \\
  \midrule
  约束方向 & 视频约束音频，单向 & 每层双向交叉注意力 & 每层双向交叉注意力 \\
  分支结构 & 两个独立模型 & 两个结构相同的分支，各约50亿参数 & 宽度可以不同的两个分支 \\
  时间对应 & 共享事件记录与时间单位 & 音频旋转位置编码按$31/157$缩放 & 跨模态连接使用时间位置 \\
  训练 & 各自训练，可选专用模型 & 音频分支先预训练，融合时冻结前馈网络 & 两路速度共同监督 \\
  推理调节 & 可单独重做任一阶段 & 共同求解，标准无分类器引导 & 文字与模态引导分别加权 \\
  主要局限 & 画面无法考虑声音，口型难对 & 约五秒片段，需配对数据 & 每步三次前向 \\
  \bottomrule
 \end{tabularx}
 \caption{音视频协同生成方式的比较。三者面对同一文字请求，差别在于两种媒体何时、以何种方向互相约束。}
 \label{tab:vision-joint-compare}
\end{table*}
\FloatBarrier

\section{多视角与三维内容生成}
\label{sec:vision-multigen-3d}

空间内容生成的输出是一个可以从不同方向观察的对象或场景，形式可以是辐射场、网格或高斯集合。难点是多视角一致：各个方向的画面必须由同一个三维表面产生。二维生成模型从海量图像学到了很强的外观先验，却没有被要求让不同视角互相一致；而带三维真值的数据远少于图像数据。另一难点是区分证据：哪些表面由观测支持，哪些由先验补出。

本节用到两组已有材料，这里只回顾结论。第一组是\ref{chap:vision-image-formation}章的相机几何：针孔相机用内参和外参把三维点投影到像素；\ref{sec:vision-formation-two-view}节说明同一三维点在两个视图中的投影受对极约束，相机已知时可以三角化出点的位置；\ref{sec:vision-formation-multi-view}节的运动恢复结构与光束法平差在多幅图中同时估计相机与三维点。这些结论告诉我们，“多视角一致”有一个可检查的几何含义：同一表面点在各视图中的投影必须落在相机几何规定的位置上。第二组是\ref{sec:vision-generation-novel-view}节的场景表示：NeRF用网络表示空间中每一点的密度和颜色，沿相机射线做体渲染得到像素（\ref{sec:vision-generation-nerf}节）；3D Gaussian Splatting用一组带中心、协方差、不透明度和颜色的三维高斯表示场景，投影到图像平面后按深度做透明度合成（\ref{sec:vision-generation-gaussians}节）。两种表示都是可微的，所以可以用渲染图与目标图之间的误差反向更新三维参数。本节的三种方法都建立在这两组材料上，差别在于目标图从哪里来：来自二维生成模型的判断、来自二维生成模型一次生成的多视角图，还是来自真实观测。

本节按输入能提供多少观测划分小节，每种条件对应一种代表方法，它们对“怎样把二维先验和三维一致性结合”给出不同回答。没有参考观测时，DreamFusion保留一个三维表示，反复从随机视角渲染，用冻结的二维扩散模型给每幅渲染图提供修改方向，一致性由共享的三维表示保证。只有单幅参考时，Wonder3D先让二维扩散模型一次生成多个视角的颜色和法线图，在生成阶段就用跨视角注意力协调一致性，再把这些图融合成表面。有多幅或连续观测时，4D Gaussian Splatting不再需要生成先验，直接用观测拟合一个随时间变化的高斯场景。三者都可能输出渲染图，但展示用的图像、可编辑网格和物理模拟模型有不同的验收要求。

\subsection{无参考观测下的生成}
\label{sec:vision-multigen-no-reference}
\label{sec:vision-multigen-dreamfusion}

文字“陶瓷玩具狗”没有给出目标对象的真实样子。若用二维生成模型分别画出正面、侧面和背面，每幅都可能画出一张脸，或者耳朵形状各不相同，因为每次生成都独立进行。训练一个直接输出三维的生成模型又缺少足够的带文字三维数据。DreamFusion的主要想法是：为每条文字请求从随机初始化的三维辐射场出发，反复渲染随机相机视图，用冻结的二维扩散模型判断渲染图应向哪个方向改，再把这个方向传回三维参数\scite{vision-dreamfusion2023}。不同视角共享同一个辐射场，一致性由表示本身保证；外观知识则全部来自二维模型，不需要任何三维训练数据。它的\term{分数蒸馏采样}（score distillation sampling，SDS）就是把二维先验偏好的变化传回三维参数的计算方法。

三维表示基于mip-NeRF 360的辐射场：一个MLP读取空间位置，输出体密度和反照率，即材料本身的颜色倾向，不含光照。这个MLP由5个残差块组成，隐藏宽度128，使用Swish激活，块间加层归一化；输入位置先经最高频率为$2^{7}$的8级正弦位置编码，并按mip-NeRF的做法对每个射线区间积分。密度经指数函数、反照率经sigmoid输出，密度上还加一个以原点为中心的高斯形偏置，使优化初期物体集中在场景中央。另有一个小MLP读取射线方向，输出环境背景，按累积不透明度与前景合成。渲染时先由密度的梯度求表面法线，再用随机采样的点光源对反照率做着色，体渲染见\ref{sec:vision-generation-nerf}节。着色迫使几何本身显出凹凸，若模型只把凹凸“画”在颜色上，换一个光照方向就会露馅。二维先验是Imagen的$64\times64$基础模型，以T5-XXL文字嵌入为条件，不使用其超分辨率级联，参数在整个过程中固定。这里设场景参数为$\bm\theta$，相机和光照条件为$\pi$，渲染结果展平为$\bm x=g_{\bm\theta}(\pi)\in\mathbb R^d$。噪声水平用$\ell$表示，加入噪声得到$\bm z_\ell=\alpha_\ell\bm x+\sigma_\ell\bm\epsilon$，其中$\bm\epsilon\sim\mathcal N(\bm0,\bm I)$。冻结模型在文字$c$下预测噪声$\widehat{\bm\epsilon}_\phi$。

SDS给出的参数梯度估计为
\begin{equation}
 \bm g_{\mathrm{SDS}}
 =w(\ell)\bm J_g^\top
     [\widehat{\bm\epsilon}_\phi(\bm z_\ell,\ell,c)
       -\bm\epsilon],
 \qquad
 \bm J_g=\frac{\partial g_{\bm\theta}(\pi)}
                 {\partial\bm\theta}\text{。}
 \label{eq:vision-sds-gradient}
\end{equation}
若参数数目为$p$，$\bm J_g\in\mathbb R^{d\times p}$，噪声差为$d$维，结果是$p$维参数梯度。白话地说，二维模型看到加噪后的渲染图，预测“其中哪些是噪声”；预测噪声与真正加入的噪声之差，指出渲染图中不像“陶瓷玩具狗”的成分，渲染Jacobian再把这个像素方向换算成参数方向。优化器沿其负方向更新场景。二维模型参数$\phi$保持固定，计算中也不对噪声预测网络的输入Jacobian反向传播。

后一限制不是普通噪声平方误差求导后的自然结果。若直接最小化
$\lVert\widehat{\bm\epsilon}_\phi(\bm z_\ell,\ell,c)-\bm\epsilon\rVert_2^2$，
链式法则还会出现噪声网络对$\bm z_\ell$的导数，它需要穿过整个大网络反传，计算昂贵且在小噪声水平下数值不稳定。SDS通过概率分布之间的目标获得另一种方向，可以由下面的简化推导看清。

固定$\pi$和$\ell$，令$q_{\bm\theta}$为均值$\alpha_\ell g_{\bm\theta}(\pi)$、方差$\sigma_\ell^2\bm I$的高斯分布，$p_\phi(\bm z_\ell\mid c)$为扩散先验在该噪声水平的条件分布。我们希望渲染图加噪后的分布落在先验分布的高概率区域，即减小$\operatorname{KL}(q_{\bm\theta}\|p_\phi)$。因为$q_{\bm\theta}$的方差不变，熵不随均值改变，KL散度对$\bm\theta$的梯度只来自交叉熵项$-\mathbb E_{q}[\log p_\phi]$。利用重参数化$\bm z_\ell=\alpha_\ell g_{\bm\theta}(\pi)+\sigma_\ell\bm\epsilon$把期望移到固定分布的$\bm\epsilon$上，再对$\bm\theta$求导，得到
\[
 \nabla_{\bm\theta}\operatorname{KL}(q_{\bm\theta}\|p_\phi)
 =-\alpha_\ell\bm J_g^\top
     \mathbb E_{\bm\epsilon}
       [\nabla_{\bm z_\ell}\log p_\phi(\bm z_\ell\mid c)]\text{。}
\]
用噪声网络提供的分数近似
$\nabla_{\bm z_\ell}\log p_\phi\approx
-\widehat{\bm\epsilon}_\phi/\sigma_\ell$，
再乘以权重$w(\ell)\sigma_\ell/\alpha_\ell$，得到
$w(\ell)\bm J_g^\top\mathbb E[\widehat{\bm\epsilon}_\phi]$。
固定相机和噪声水平时，$\bm J_g$不依赖抽样噪声，而$\mathbb E[\bm\epsilon]=\bm0$，因此可减去$\bm\epsilon$作为零均值控制变量，形成式\eqref{eq:vision-sds-gradient}的估计。减去$\bm\epsilon$不改变期望，但当网络预测接近真实噪声时，单次估计的方差明显变小。

这一推导要求所用噪声水平的$\alpha_\ell,\sigma_\ell$非零，并把学习到的分数当作分布分数的近似。实际引导和网络误差会影响方向，不能因此声称每步都精确降低真实三维分布的KL散度。实现可对$w(\ell)(\widehat{\bm\epsilon}_\phi-\bm\epsilon)$停止梯度，再让它与渲染图像作内积求导；直接对噪声差的平方求导会得到不同计算。

用一个局部线性教学渲染器检查方向。设$\bm\theta=(a,b)^\top$，两个像素分量为$g_{\bm\theta}=(a,a+b)^\top$，则
$\bm J_g=\left(\begin{smallmatrix}1&0\\1&1\end{smallmatrix}\right)$。
若权重为1、噪声差为$(0.2,-0.1)^\top$，参数梯度为$(0.1,-0.1)^\top$。从$(0.5,0.2)^\top$以步长0.1更新，得到$(0.49,0.21)^\top$，渲染变为$(0.49,0.70)^\top$。第二像素同时受$a,b$影响，两者变化抵消，说明像素方向不能不经渲染Jacobian直接当作参数方向。图\ref{fig:vision-sds-path}标出了梯度经过和停止的位置。

\begin{figure*}[tbp]
  \centering
  \begin{tikzpicture}[x=1cm,y=1cm,font=\figurefont\small,
  block/.style={draw=FigureModel,fill=FigureModel!10,line width=0.8pt,
    minimum height=10mm,minimum width=23mm,align=center},
  input/.style={block,draw=FigureInput,fill=FigureInput!10},
  arr/.style={-{Stealth[length=2mm]},BookInk,line width=0.8pt,
    rounded corners=1.2mm},
  grad/.style={-{Stealth[length=2mm]},FigureLoss,dashed,line width=1pt,
    rounded corners=1.2mm}]
  \node[input] (theta) at (1.2,2) {三维参数\\$\bm\theta$};
  \node[block] (render) at (4.5,2) {渲染器\\$g_{\bm\theta}(\pi)$};
  \node[block] (noise) at (7.8,2) {加噪\\$\alpha_\ell\bm x+\sigma_\ell\bm\epsilon$};
  \node[block] (prior) at (11.1,2) {冻结二维先验\\$\widehat{\bm\epsilon}_\phi$};
  \node[block] (res) at (14.6,2) {噪声差\\$\widehat{\bm\epsilon}_\phi-\bm\epsilon$};
  \draw[arr] (theta) -- (render);
  \draw[arr] (render) -- node[above] {$\bm x$} (noise);
  \draw[arr] (noise) -- (prior);
  \draw[arr] (prior) -- (res);
  \node (camera) at (4.5,3.5) {相机与光照$\pi$};
  \node (eps) at (7.8,3.5) {噪声$\bm\epsilon$};
  \node (text) at (11.1,3.5) {文字$c$};
  \draw[arr] (camera) -- (render);
  \draw[arr] (eps) -- (noise);
  \draw[arr] (eps.north) -- (7.8,4.25) -- (14.6,4.25) -- (res.north);
  \draw[arr] (text) -- (prior);
  \node[block,draw=FigureLoss,fill=FigureLoss!8,minimum width=35mm]
    (jac) at (4.5,-0.15) {通过渲染器求导\\$w(\ell)\bm J_g^\top(\widehat{\bm\epsilon}_\phi-\bm\epsilon)$};
  \draw[grad] (res.south) -- (14.6,-0.15) -- node[above] {噪声差停止梯度} (jac.east);
  \draw[grad] (jac.west) -- (1.2,-0.15) -- (theta.south);
  \node[anchor=west] at (7.1,-0.8) {不反传噪声网络的输入Jacobian};
\end{tikzpicture}

  \caption{SDS把冻结二维先验的噪声差传回三维参数。实线为渲染、加噪与预测，虚线为通过渲染器的参数梯度。噪声差在求参数梯度时视为常量，不反传冻结噪声网络的输入Jacobian；每次重新抽取相机，使多个方向约束同一个场景。}
  \label{fig:vision-sds-path}
\end{figure*}
\FloatBarrier

完整的优化过程对每条文字请求单独进行，没有跨请求的训练阶段。每次迭代随机抽取相机位置、焦距和点光源，渲染$64\times64$图像，从$[0.02,0.98]$中均匀抽取噪声时刻，计算式\eqref{eq:vision-sds-gradient}并更新辐射场参数；论文取权重$w(\ell)=\sigma_\ell^2$，并报告均匀权重效果相近。相机仰角在$-10^\circ$到$90^\circ$之间、方位角在$0^\circ$到$360^\circ$之间抽取，相机到中心的距离在1到1.5之间。论文共优化15000步，用Distributed Shampoo优化器，在4块TPUv4芯片上约需1.5小时。除SDS外，还加入表面朝向正则项，惩罚背向相机的法线，以及控制累积不透明度、抑制半透明浮游物的正则项。论文使用很大的无分类器引导权重100，这使渲染结果更贴近文字，但也使颜色偏饱和、几何偏平滑。模型还按相机方位角在文字后附加“正面视角”“侧面视角”“背面视角”等描述，减少把背面也拉向正面外观的冲突。

优化完成后，辐射场可以从不同相机渲染；若要交付网格，还须提取表面并整理材质。多个正面、几何过平滑、细结构缺失或依赖视角的伪外观，都可能在环绕观察中出现。多个正面的根源在于二维先验对每幅渲染图独立判断“像不像玩具狗”，而训练图像中正面照最多，视角描述只能部分抑制这一偏好。每条请求都要重新优化，一次生成按论文设置约需1.5小时。可渲染性没有证明表面可直接编辑，生成的背面也不是任何真实玩具的背面。

\subsection{单幅参考下的生成}
\label{sec:vision-multigen-single-reference}
\label{sec:vision-multigen-wonder3d}

单幅照片增加了可见轮廓、颜色和主体身份的约束，但背面和遮挡区域仍未知。创作前可约定正面保持原耳朵、眼睛与花纹，允许设计背面纹理。这样才能把偏离参考的错误，与合理补全但无法验证的内容区分开。

DreamFusion式的逐请求优化在这里有两个不足：每个对象都要优化很久，而且二维先验对每幅渲染图独立打分，多视角一致性要靠长时间优化慢慢形成。另一种做法是先让二维扩散模型一次生成多个固定视角的颜色图，再从这些图重建三维。但只生成颜色图，重建时仍可能把纹理变化解释成凹凸，因为颜色不直接说明表面朝向。Wonder3D的主要想法是同时生成多视角颜色图和\term{法线图}（normal map），让外观与表面朝向在生成阶段就互相约束，再用一个快速的表面优化把它们融合\scite{vision-wonder3d2024}。法线图的每个有效像素给出该处表面的朝向向量，它不直接给出深度或三维点位置。

生成部分以单幅$256\times256$参考图为条件，输出前、后、左、右、右前、左前六个固定视角的颜色图和法线图，共12幅。去噪网络由Stable Diffusion Image Variations模型初始化，它本来就以图像为条件生成图像。输出域由一个一维的域标记说明当前生成的是颜色还是法线，该标记经位置编码后与生成时刻嵌入拼接，进入U-Net。原来的自注意力被扩展为跨视角注意力：一个视角的查询可以读取全部六个视角的键和值，使不同相机下的表示互相协调。每个Transformer块中还在交叉注意力之前加入跨域注意力，把同一视角的颜色与法线表示的键和值合并。仅有域标记能告诉网络输出哪类数据，却不能单独保证两类结果描述同一个表面，跨域注意力正是为此加入的。

两种注意力的差别在于键和值从哪里取。把12幅图在U-Net某一层的特征记为$\bm H_{k,u}$，其中$k=1,\dots,6$为视角，$u\in\{\mathrm{rgb},\mathrm{n}\}$为颜色或法线域，每幅图有$m$个空间位置。多视角注意力中，视角$k$、域$u$的查询读取同一域全部六个视角的键和值：
\begin{equation}
 \operatorname{MVAttn}(\bm H_{k,u})
 =\operatorname{Attn}\bigl(\bm Q_{k,u},
   [\bm K_{1,u};\dots;\bm K_{6,u}],
   [\bm V_{1,u};\dots;\bm V_{6,u}]\bigr),
 \label{eq:vision-multigen-wonder-mv}
\end{equation}
键值序列长$6m$。跨域注意力中，查询读取同一视角两个域的键和值：
\begin{equation}
 \operatorname{CDAttn}(\bm H_{k,u})
 =\operatorname{Attn}\bigl(\bm Q_{k,u},
   [\bm K_{k,\mathrm{rgb}};\bm K_{k,\mathrm n}],
   [\bm V_{k,\mathrm{rgb}};\bm V_{k,\mathrm n}]\bigr),
 \label{eq:vision-multigen-wonder-cd}
\end{equation}
键值序列长$2m$。前者没有显式使用对极约束，相机之间的对应由注意力权重从数据中学出；后者让同一视角的颜色与法线互相参照。以$32\times32$的特征层为例，$m=1024$，多视角注意力的键值长度为6144，普通自注意力只有1024，代价随之增加。论文说明跨域注意力放在U-Net每个Transformer块的交叉注意力之前，多视角注意力由原自注意力扩展而来；注意力头数等其余配置沿用初始化模型，论文未单列。

训练使用Objaverse中经过清理的约三万个三维对象。对每个对象从固定视角渲染颜色图和法线图，以某一视图作为输入参考，其余视图和法线作为目标，用标准的去噪扩散目标微调网络。论文报告训练3万步，总批量512，在8块A800 GPU上约需三天；训练沿用初始化模型的噪声预测目标与优化器设置。采样时从12幅图的噪声共同生成结果，未观测的背面由生成先验补充。六幅看起来各自合理的图，仍可能在细耳朵、尾巴和遮挡位置互相矛盾。

形成几何时，方法优化一个\term{有符号距离场}（signed distance field，SDF）：空间点到表面的距离带有内外符号，零值集合定义表面。若约定物体外部为正，在可微且梯度非零处，
$\bm n(\bm p)=\nabla f(\bm p)/\lVert\nabla f(\bm p)\rVert_2$
给出向外单位法线。法线由SDF对空间的一阶导数得到；优化涉及法线的损失时，还要对场参数继续求导，不能把“求法线”本身称为计算SDF的二阶空间导数。SDF用instant-NGP式的多分辨率哈希网格表示，以加快优化。

从各生成视图抽取像素射线，用SDF渲染颜色、法线和前景占据，再与相应生成图比较。颜色用平方误差，前景掩码用二元交叉熵，法线用下面的加权损失；梯度范数接近1的正则项保持距离场性质，稀疏与平滑项限制浮游表面和局部噪声。论文把总目标写成这六项之和，各项权重未在文中给出。这些是对生成多视图的融合，不是把六张图直接粘成六面体。

生成法线也可能错误，因此原方法依据观察方向赋权。设$\bm v_i$是从相机射入表面的单位方向，$\bm n_i$是生成的向外单位法线，$\widehat{\bm n}_i$是当前渲染法线。几何法线损失为
\begin{equation}
 L_{\mathrm n}
 =\frac{\sum_i w_i(1-\widehat{\bm n}_i^\top\bm n_i)}
         {\sum_i w_i},
 \quad
 w_i=
 \begin{cases}
  0,&\bm v_i^\top\bm n_i>\eta,\\
  \exp(|\bm v_i^\top\bm n_i|),&\bm v_i^\top\bm n_i\leq\eta,
 \end{cases}
 \label{eq:vision-wonder-normal}
\end{equation}
其中$\eta$为接近0的负阈值，且当前批次要求$\sum_iw_i>0$。相机看到的向外表面法线应大致逆着入射视线，点积应为负；接近掠射的法线获得较低权重，朝向明显矛盾的结果被排除。全被排除时不能除以零，应跳过该项或重新取样。

取两个教学射线，方向点积为$-1$和$-0.5$，法线差异$1-\widehat{\bm n}_i^\top\bm n_i$分别为0.1和0.3，阈值取$-0.01$。两项权重为$e$和$e^{0.5}$，归一化损失约0.17551，比等权平均0.2更靠近正视方向的误差。它说明赋权的含义，不证明正视法线一定正确。颜色和掩码误差还采用剔除一部分最大残差的策略，降低局部不一致的影响，但可能同时忽略真实细节，须观察最终表面。论文报告整个流程约两分钟即可得到带纹理的网格，远快于逐请求的SDS优化。

输入玩具狗正面时，可以检查输入侧轮廓是否保持，再环绕检查耳朵是否连接、纹理是否有接缝。六个固定视角之间的斜后方没有生成图约束，那里的表面只由SDF的平滑性插出。若没有真实尺度或额外视图，背面纹理与绝对尺寸都不具有测量依据。用于展示的成果也许已经足够，用于动画绑定或制造则还要检查拓扑、厚度与材质等要求。

\subsection{多幅与连续观测下的生成}
\label{sec:vision-multigen-observed-space}
\label{sec:vision-multigen-4dgs}

多幅观测能限制更多表面，前提是相机和时间关系正确。静态场景的拟合与新视角合成已在\ref{sec:vision-generation-novel-view}节展开。若玩具狗正在运动，就不能把不同时间的腿部位置全部拟合为同一静态表面；若为每个时刻单独拟合一组高斯，每个时刻只有少数相机看到，表面会严重欠约束，存储也随帧数线性增长。需要一个随时间变化、又能在时刻之间共享信息的场景表示。

Wu等人的4D Gaussian Splatting的主要想法是：保留一组所有时刻共享的基准三维高斯，用一个紧凑的形变场根据位置和时间预测每个高斯的变化，再渲染相应时刻\scite{vision-4dgs2024}。共享基准使每个时刻都能借用其他时刻的观测，形变场则只需表示变化部分。这里的“四维”指三维空间随时间变化，采用的是该文的时空平面与形变路线，不与其他同名的四维高斯参数化混用。

静态高斯的中心、旋转、尺度、不透明度和颜色，以及投影与透明度合成，见\ref{sec:vision-generation-gaussians}节。形变模块读取基准中心$\bm\mu_i=(x_i,y_i,z_i)^\top$及查询时间$\tau$，在$(x,y)$、$(x,z)$、$(y,z)$、$(x,\tau)$、$(y,\tau)$和$(z,\tau)$六种二维平面上取特征。每个平面用四个邻近网格点双线性插值，因此不用存储完整四维稠密网格。论文的平面基础分辨率为64，并在此基础上再上采样2倍和4倍形成多分辨率（D-NeRF合成数据只用2倍一级）。若每个方向都取64个网格点，四维稠密网格有$64^4\approx1.7\times10^7$个格点，六个平面合计只有$6\times64^2=24576$个；这里假设时间方向也取64个格点只是为了比较，论文没有单独给出时间方向的分辨率，也没有给出平面特征的通道数。

在同一分辨率下，六个平面特征逐元素相乘；不同分辨率的结果拼接，再由一个小型MLP $\phi_{\mathrm d}$合并为形变特征$\bm f_{\mathrm d}$，最后由三个分开的MLP头分别输出位置、旋转和尺度的变化。设$\mathcal P$为上述六组坐标对，$r$标记分辨率，可写
\begin{equation}
 \begin{aligned}
 \bm h_i(\tau)
 &=\mathop{\operatorname{concat}}_r
   \left[
    \mathop{\bigodot}_{(a,b)\in\mathcal P}
       \operatorname{interp}
        (\bm F_{a,b,r},a_i,b_i)
   \right],\\
 \bm f_{\mathrm d}&=\phi_{\mathrm d}(\bm h_i(\tau)),\\
 \Delta\bm\mu_i(\tau)&=\phi_{\mathrm x}(\bm f_{\mathrm d}),\quad
 \Delta\bm r_i(\tau)=\phi_{\mathrm r}(\bm f_{\mathrm d}),\quad
 \Delta\bm s_i(\tau)=\phi_{\mathrm s}(\bm f_{\mathrm d})\text{。}
 \end{aligned}
 \label{eq:vision-4d-planes}
\end{equation}
当坐标对含$\tau$时，式中的相应坐标取查询时间；$\bm F_{a,b,r}$表示带特征通道的网格数组，$\Delta\bm r_i$与$\Delta\bm s_i$是旋转和尺度参数的变化量。论文没有给出$\phi_{\mathrm d}$与三个头的层数和宽度，只在HyperNeRF数据上说明把解码器宽度改为128。相邻位置会共享网格参数，使附近高斯的形变具有共同表示，但不会强制它们属于同一个刚体。

中心按$\bm\mu_i(\tau)=\bm\mu_i+\Delta\bm\mu_i(\tau)$更新，旋转和尺度也在相应参数化中加入预测变化。渲染前还须把旋转参数恢复为合法旋转、尺度恢复为正值，不能把任意残差直接当作正定协方差。原路线的高斯不透明度与颜色参数仍来自基准集合，查询相机时再进行投影与透明度合成。

教学上，一个中心原为$(0,0,2)^\top$，在两个时刻预测的位移分别为$(0,0,0)^\top$和$(0.2,0,0)^\top$。若固定针孔相机焦距为100像素、主点横坐标为0，投影横坐标从0变为$100\times0.2/2=10$像素。纵深、旋转和尺度变化还会影响投影足迹，不能只把二维图中的点平移而忽略高斯形状。

训练数据是带相机参数和时间戳的多视角视频或单相机连续视频。训练先不使用形变场，用3000次迭代优化静态高斯获得初始场景，再联合学习基准表示与形变场；D-NeRF合成数据共训练20000次迭代，Neu3D多相机数据为14000次。每次选择观测相机及时间，渲染当时的高斯，以对应图像的L1像素误差作为监督，并对时空平面施加基于网格的全变差正则，使形变在空间和时间上连续。论文报告在单块RTX 3090上，D-NeRF合成场景以$800\times800$分辨率渲染可达每秒82帧，存储约18 MB。使用时输入任意查询相机和时间，生成该位置的视图；基准高斯、平面特征和小型网络共同构成需要保存的场景。

验收应把已观测时间、未观测视角和观测间时间分开。训练相机下拟合良好，可能仍在背面或快速运动时失败。单相机连续观测还把视角与时间交织在一起，某些表面变化难以唯一解释。高斯有基准编号也不等于它们具有真实物体的持久身份，形变网络更不是含质量、力和碰撞约束的物理模拟器。

三种方法的组件配置列在表\ref{tab:vision-multigen-3d-config}中。三者都不是单一的Transformer，表的列因此按组件而不是按层数、宽度排列。

\begin{table*}[tbp]
 \centering
 \small
 \begin{tabularx}{\linewidth}{@{}p{22mm}XXX@{}}
  \toprule
  组件 & DreamFusion & Wonder3D & 4D Gaussian Splatting \\
  \midrule
  三维表示 & mip-NeRF 360式MLP：5个残差块、宽128、8级积分位置编码；另有背景MLP & instant-NGP式哈希网格SDF & 基准三维高斯集合 \\
  生成或形变网络 & 冻结的Imagen $64\times64$基础模型，T5-XXL文字条件 & 由Stable Diffusion Image Variations初始化的U-Net，加多视角与跨域注意力 & 六平面特征（基础分辨率64，再上采样2倍、4倍）、合并MLP与三个形变头 \\
  分辨率 & 渲染$64\times64$ & 参考与生成图$256\times256$，6视角$\times$2域 & 随数据集，例如D-NeRF为$800\times800$ \\
  训练或优化 & 每条请求15000步 & 生成网络3万步，批量512；表面融合每个对象单独优化 & 静态预热3000次，再联合优化至14000至20000次 \\
  报告耗时 & 4块TPUv4约1.5小时 & 训练8块A800约三天；单对象约两分钟 & 单块RTX 3090，渲染每秒30至82帧 \\
  \bottomrule
 \end{tabularx}
 \caption{三种空间内容生成方法的组件配置。数值取自各论文；注意力头数、平面特征通道数、形变MLP层数等论文未给出的项不列。}
 \label{tab:vision-multigen-3d-config}
\end{table*}
\FloatBarrier

三种条件下的方法比较见表\ref{tab:vision-3d-compare}。从DreamFusion到Wonder3D，改变的是一致性在哪里形成：前者在优化中靠共享表示逐步形成，后者在生成时就靠跨视角注意力协调，换来了从小时到分钟的速度，代价是需要三维资产渲染数据训练，且只覆盖六个固定视角。4D Gaussian Splatting去掉了生成先验，结果有观测支持，但它只能重现被拍到的内容，不能为未拍到的背面创作外观。

\begin{table*}[tbp]
 \centering
 \small
 \begin{tabularx}{\linewidth}{@{}p{22mm}XXX@{}}
  \toprule
  比较项 & DreamFusion & Wonder3D & 4D Gaussian Splatting \\
  \midrule
  输入条件 & 文字 & 单幅参考图 & 多视角或连续观测及相机 \\
  三维表示 & 辐射场（密度与反照率） & 有符号距离场 & 基准高斯加时空平面形变场 \\
  一致性来源 & 共享表示，随机视角反复优化 & 跨视角与跨域注意力，再融合表面 & 观测像素与共享基准 \\
  先验与训练 & 冻结二维扩散模型，无三维数据 & 用三维资产渲染微调二维扩散模型 & 无生成先验，逐场景拟合 \\
  单次耗时 & 逐请求优化，以小时计 & 约两分钟 & 逐场景优化 \\
  未见区域依据 & 全部来自先验 & 背面来自先验 & 无观测处缺乏依据 \\
  \bottomrule
 \end{tabularx}
 \caption{三种观测条件下空间内容生成的比较。耗时依据各论文报告，硬件不同，只说明量级。}
 \label{tab:vision-3d-compare}
\end{table*}
\FloatBarrier

\section{视觉条件下的文字与语音响应}
\label{sec:vision-multigen-response}
\label{sec:vision-multigen-omni-response}

当用户询问“哪个物体先碰到桌面”，要生成的是基于视觉证据的回答，以及说出回答的声音。答案是否有依据属于\ref{chap:vision-multimodal-understanding}章的问题；这里关注文字内容如何转成语音，以及用户何时能听到它。输入是视频、声音和问题，输出是文字与语音两路回答；难点在于两路内容必须一致，而且语音要尽早开始播放，又不能因为抢先开口而说出证据尚不支持的结论。这与给视频中的物体配音效具有不同目标，与前几节的媒体生成相比，生成内容由理解结果决定。

Qwen2.5-Omni的结构与训练见\ref{sec:vision-models-qwen-omni}节：Thinker读取视觉、声音和问题，生成文字及其隐表示；Talker接收这些隐表示和已采样的文字词元，自回归地产生语音编码；流式解码器把语音编码分块转换为梅尔频谱，再由声码器恢复波形\scite{vision-qwenomni2025}。Talker同时接收文字词元，帮助区分语义接近但发音不同的词，而不只依据模糊的语义向量猜测音节。

使用时一般不微调模型，而是组织输入和设置输出方式。视频按实际时间与声音一起送入，模型用对齐的时间位置表示两者；问题以文字或语音给出。输出端要选择是否需要语音，只需文字时可以不运行Talker，减少计算。语音编码需要解码才能播放：原方法把相邻编码分成块，解码当前块时读取有限的前后上下文，报告中的设置包含前面两块、当前块和后面一块。因此“流式”并不意味着拿到第一个文字词元就能零等待地发声；前看一块本身就需要一定积累。块越大，每块需要等待的生成越多，首次发声越晚；块越小，块边界越多，声音可能出现不连续。

响应延迟至少要区分首次发声和整段完成。设用户输入就绪时刻为0，第$k$块可播放音频准备好于$a_k$，播放时长为$d_k$，实际开始播放的时刻为$b_k$。忽略额外播放器缓冲时，
\begin{equation}
  b_1=a_1,\qquad
  b_k=\max(a_k,b_{k-1}+d_{k-1})\text{，}\quad k>1\text{。}
  \label{eq:vision-stream-playback}
\end{equation}
若音频尚未准备好，就必须等到$a_k$；若准备得很早，也要等前一块播完。这个排程关系与模型的具体参数无关，可以用来检查推理和播放是否衔接。

例如三块音频分别在0.8、1.2、1.7秒就绪，每块长0.5秒，则开始播放时刻为0.8、1.3、1.8秒，整段在2.3秒结束，中间无空隙。若第二块直到1.5秒才就绪，开始时刻变成0.8、1.5、2.0秒，出现0.2秒停顿。全部生成后再播放，在前一设置中要等到1.7秒才开始，并于3.2秒结束。以上均为教学排程，不是Qwen2.5-Omni的性能实测。首次发声时间包含输入编码、文字生成、语音编码预测及首块波形解码；视频越长，输入编码越久，首次发声越晚。

视觉问答还需决定是否可以提前说出某个结论。若视频中第二次接触尚未观察到，提前回答“花瓶先落桌”可能建立在猜测上。已经播放的词句无法像屏幕文字那样悄然撤回，纠正需要明确说出来。流式系统应在证据已经足够时输出相应内容，并把后来获得的信息与已说出的内容保持一致。

常见失败包括：语音漏读否定词，使“没有碰到”变成“碰到”；把“左边”说成“右边”；在数字中间断块造成停顿；以及文字回答正确而语音转写不同。验收时同时查看视觉证据、显示文字和语音转写，自然度、可懂度、首次发声时间和最长停顿分别记录；更早开口不能补偿答错。该模型生成文字与语音回答，不承担视频内容生成。

\section{成果验收}
\label{sec:vision-multigen-evaluation}

前几节在每种成果之后都留下了核验对象：视频的身份与运动、编辑的保留范围、声轨的起声时刻、联合音画的共同事件、三维内容的观测支持、语音回答的时序。本节把它们组织成交付前的验收。媒体成果既要独立可用，也要满足请求中规定的共同关系，还要说明为得到它花了多少代价。三者分别回答“单项好不好”“关系对不对”“代价多大”，按维度划分小节；把它们保留为分项结果，才能判断失败发生在哪里。自动特征分数、人类观看试听和具体任务检查承担不同职责，下文在各维度中分别说明；\ref{sec:vision-multigen-challenges}节定义的分布距离、逐片段指标和同步偏移，在这里按它们回答的问题归入相应维度。

\subsection{单项质量的核验}
\label{sec:vision-multigen-quality}

视频可检查模糊、闪烁、运动和身份变化；声音可检查失真、声源、语言内容和可懂度；空间资产可检查表面、纹理、拓扑和编辑用途。指标应作用于相应交付对象。图像质量分数不包含声音，可旋转展示的辐射场也不自动满足网格编辑要求。

视频的单项质量最难用一个数概括，两个基准给出不同的分解方式。VBench把质量拆成若干相互独立的维度，每个维度用一种专门工具测量画面层面的性质；VBench-2.0在此之上转向事件层面，问生成内容是否符合物理和常识。前者回答“看起来是否稳定、清晰”，后者回答“发生的事是否合理”。

\subsubsection{VBench}
\label{sec:vision-multigen-vbench}

单个总分，例如式\eqref{eq:vision-multigen-fvd}的FVD这类生成视频与真实视频的特征分布距离，会把闪烁、身份变化和文字不符混在一起，无法说明模型在哪方面失败。VBench把视频生成评价拆成16个层次化、相互解耦的维度，包括主体与背景一致性、时间闪烁、运动平滑度、动态程度、画面质量及条件语义等，再为每个维度选择专门的分析工具。其用途是定位差异，避免一个总分遮住主要失败\scite{vision-vbench2024}

十六个维度分为两大类。视频质量一类不看提示内容，又分时间质量与逐帧质量；视频与条件一致性一类检查提示中的语义和风格是否出现，最后另有一个整体一致性维度。表\ref{tab:vision-multigen-vbench-dims}列出每个维度所用的工具和计分方式。主体一致性已在式\eqref{eq:vision-multigen-subject-consistency}中给出，背景一致性用同一公式，只是把DINO特征换成CLIP图像特征。时间闪烁在提示要求静止场景的视频上计算相邻帧的平均绝对像素差$S$，再换算为$(255-S)/255$，越接近1越稳定。运动平滑度去掉奇数帧，用AMT插帧模型从偶数帧重新插出它们，再与原奇数帧比较平均绝对误差，误差小说明运动符合插帧模型所学的平滑先验。动态程度用RAFT估计光流，取每段视频中最大的5\%光流幅值的平均来判断视频是否在动，分数是判为非静止的视频比例。这些工具输出的是代理测量，例如特征相似度高不能证明每根手指都保持正确。

\begin{table*}[tbp]
 \centering
 \small
 \begin{tabularx}{\linewidth}{@{}p{22mm}p{26mm}p{34mm}X@{}}
  \toprule
  类别 & 维度 & 工具 & 计分方式 \\
  \midrule
  时间质量 & 主体一致性 & DINO & 式\eqref{eq:vision-multigen-subject-consistency}的帧间余弦相似度 \\
   & 背景一致性 & CLIP图像编码器 & 同一公式，换用CLIP特征 \\
   & 时间闪烁 & 像素差 & 静止场景相邻帧的平均绝对差，换算为$(255-S)/255$ \\
   & 运动平滑度 & AMT插帧模型 & 由偶数帧重插奇数帧，与原帧比较 \\
   & 动态程度 & RAFT光流 & 判为非静止的视频比例 \\
  逐帧质量 & 美学质量 & LAION美学预测器 & 0至10的逐帧评分线性归一化后平均 \\
   & 成像质量 & 在SPAQ上训练的MUSIQ & 逐帧分数除以100后平均 \\
  语义 & 对象类别 & GRiT检测 & 检出指定类别的帧比例 \\
   & 多个对象 & GRiT检测 & 所有指定对象同时出现的帧比例 \\
   & 人体动作 & 在Kinetics-400上训练的UMT & 均匀取16帧分类，检查目标动作是否在高置信候选中 \\
   & 颜色 & GRiT描述 & 含对象且有颜色描述的帧中，颜色正确的比例 \\
   & 空间关系 & 规则判定，沿用T2I-CompBench & 对象对的左右、上下关系正确率 \\
   & 场景 & Tag2Text描述 & 生成的场景描述是否与提示一致 \\
  风格 & 外观风格 & CLIP & 帧与风格描述的特征相似度 \\
   & 时间风格 & ViCLIP & 视频与镜头运动等时间风格描述的相似度 \\
  整体 & 整体一致性 & ViCLIP & 一般提示下的视频文字相似度 \\
  \bottomrule
 \end{tabularx}
 \caption{VBench十六个维度的评价工具。每个维度配有约100条提示的专用提示集；分数都是工具给出的代理测量，工具的误判计入评价误差。}
 \label{tab:vision-multigen-vbench-dims}
\end{table*}
\FloatBarrier

评价流程从固定提示集合开始，每个维度有自己约100条提示的子集；为每条提示保留规定数量的生成样本，用同一版本的分项脚本处理，再与相应的人类偏好标注比较。VBench用人类标注检查自动测量的对应关系，但某个基准中的对应不保证在所有新场景中成立。工具对生成缺陷的误判也应保留为评价误差。

一个几乎不动的玩具狗片段可能在外观一致性上很好，却没有完成“绕花瓶走一圈”。另一个片段运动丰富，却在转身后变成另一只狗。二者需要分别报告动作完成与身份保持，不能把低动态当成稳定性的无条件优点。模型间比较还须固定帧率、长度、样本数和生成预算，否则评分差异可能来自不同制作条件。

\Needspace{65mm}
\subsubsection{VBench-2.0}
\label{sec:vision-generation-physics-check}
\label{sec:vision-multigen-vbench2}

VBench的维度主要测量画面层面的性质。一个片段可以每帧清晰、运动平滑，却让杯子穿过桌面或让玩具狗凭空多出一条腿。这类错误不在画面层面，而在事件层面。VBench-2.0的2025年v2把评价目标改为内在的真实性，按人体表现、可控性、创造性、物理与常识五个方面组织测试，每个方面再细分为若干能力\scite{vision-vbench22025}。评价工具组合了两类：通用的视觉语言模型和大语言模型，用来回答关于视频内容的问题；以及专项检测器，例如为生成视频设计的异常检测方法。

在物理状态变化中，评价先根据提示形成具体的可见行为描述，再用多个问题核查生成视频。部分流程先排除初始状态不符合要求的片段，使状态变化判断对应原题。解释这种分数时，应同时报告初始条件的通过情况；只看过滤后的得分，会遗漏模型连题目起点都没有生成正确的失败。

对“玩具狗推倒花瓶”，教学核验可以分别问是否接触、花瓶是否倾斜、倾倒后是否保留同一个物体。仅有抬腿动作或一声碰撞都不足以回答这些问题。常识评价中的异常合并、突然复制与消失也需要跨帧观察，逐帧都像真图不表示帧间关系合理。

该基准的多视角一致性使用特征匹配等代理证据，因为生成视频通常没有对应真实三维模型。它不能直接测出真实几何误差。测试分数提高支持的是相应题目和评价流程下的改善，不是已经具备普遍物理模拟能力；视觉语言模型作为裁判也会出错，自动判定仍须与可见事实及人类检查对照。

两个基准的关系可以用同一组玩具狗视频说明。VBench能发现“转身后变成另一只狗”和“几乎不动”，因为这两者在特征相似度和光流上有明显信号；VBench-2.0能发现“玩具狗碰到花瓶，花瓶却没有倒，下一帧又出现在桌下”，因为它针对事件提问。VBench-2.0改变了评价工具的类型，从专用特征提取器换成以问题驱动的视觉语言模型，换来了对事件和常识的覆盖，代价是裁判模型本身的错误更难预测，分数的可重复性依赖裁判模型的版本。两者都只评价视频，不覆盖声轨、三维资产和语音回答。

\subsection{条件与跨输出关系的核验}
\label{sec:vision-multigen-relations}

单项质量合格，不等于成果满足请求。验收先把请求中的对象与关系落到同一份记录上。参考玩具狗的身份、杯子的材质、碰撞次数、声源及起声时刻相互关联，但应分别有可检查对象。表\ref{tab:vision-media-acceptance}列出贯穿示例中的核验方式，表中任务阈值需在生成前确定。

\begin{table*}[tbp]
 \centering
 \small
 \begin{tabularx}{\linewidth}{@{}p{25mm}p{44mm}X@{}}
  \toprule
  要求 & 主要证据 & 容易混淆的成功与失败 \\
  \midrule
  主体保持 & 原参考、转身前后与遮挡后画面 & 颜色相近仍可能改变耳朵、形状或身份 \\
  局部修改 & 原片段、编辑掩码、最终像素 & 指定区域改对，掩码外却有背景漂移 \\
  碰撞配声 & 接触时间、起声时间、事件次数 & 声音自然，但过早、过晚或多出一次 \\
  联合音画 & 共同事件、对白转写与口型 & 两路各自自然，却来自不同动作或声源 \\
  空间内容 & 参考覆盖、相机、表面与材质 & 新视图逼真，背面却无观测依据或不可编辑 \\
  语音回答 & 视觉依据、文字、语音与时序 & 文字正确，语音漏词，或早开口后长时间停顿 \\
  \bottomrule
 \end{tabularx}
 \caption{媒体制作的分项验收示例。证据应对应最终交付格式，掩码、相机、时间和事件来源需保存。各行回答不同问题，不能由某一项高分推断整件成果通过。}
 \label{tab:vision-media-acceptance}
\end{table*}
\FloatBarrier

跨输出关系可能随制作步骤改变。视频编辑把花瓶换成金属杯后，原陶瓷碰撞声可能已经不合适；三维资产改短了尾巴，沿用旧轮廓生成的多视角宣传图可能不再一致。应从受影响的关系出发重新核验，而不是只检查被修改的文件。

空间成果需标明哪些表面有输入观测支持，哪些由先验补出。可以在留出的真实视图上测量渲染误差；没有真实背面时，只能检查创作合理性、内部一致性及用途。把全部生成视角当作“验证视角”，会让同一个生成先验同时提供答案和检验依据。同样，用生成时使用的同步编码器来评价同步，也会让条件模型与评价模型重合，应当用独立的方法或人工检查补充。

\subsection{制作预算的核验}
\label{sec:vision-multigen-budget}

一条成功视频可能来自大量候选与多轮局部重做。比较制作能力时，应固定输入、时长、分辨率、模型版本、候选数量和编辑预算，并记录达到可用成果的总代价。只展示最佳样本会把筛选成本和失败概率都隐藏起来。

假设两个教学流程都交付一个四秒片段。甲生成4个候选，每个占用2分钟计算，再做一次1分钟的局部修改，累计计算量为9设备分钟；乙生成1个候选，占用5设备分钟。两者都通过当前验收，也不能因此说甲更省资源。若甲的4个候选并行生成，其等待时间可能短于9分钟，但设备累计时间仍是9分钟，人工观看与选择时间还要另计。本章的模型在代价结构上差别很大：DreamFusion每个对象要优化很久，Wonder3D约两分钟；LTX-2的模态引导每步需要三次前向；Wan的14B模型未经优化时单次约半小时。这些都应计入同一预算口径。

任务完成率的分母应是事先选定的请求，而不是最后留下的候选。可以记录规定预算内是否至少有一个结果通过，并同时给出每请求候选数、重做次数与失败类型。这样既描述了最终可用性，也保留模型一次生成的不稳定性。

\section{本章小结}
\label{sec:vision-multigen-summary}

开篇的请求要求玩具狗在转身后保持原样、碰撞真实发生、声音落在接触时刻，还要能从各个方向观察。本章逐步改变“哪些内容固定、哪些内容生成”的设定来回答这些要求。时间一致性方面，Lumiere在网络内部压缩时间，Wan先用自编码器压缩再用Transformer处理整段潜网格；修改与延伸则通过掩码、原片段或已知帧把“应当保留什么”写进条件。声音方面，MMAudio用高帧率同步特征决定何时发声，Movie Gen Audio用音频上下文延伸和修补长声轨；画面与声音都可变时，Ovi和LTX-2在每个生成步让两路互相读取，分阶段路线则以单向约束换取模块化。

空间方面，DreamFusion用分数蒸馏把二维先验传回共享的三维表示，Wonder3D在生成阶段就协调多视角颜色与法线，4D Gaussian Splatting用观测拟合随时间形变的高斯场景。三者的差别归结为一致性在哪里形成，以及未见区域由先验还是观测决定。视觉条件下的语音响应还把生成延伸到播放时序，文字正确、语音可懂和及时作答共同决定使用效果。

验收把这些要求拆回单项质量、条件与跨输出关系、制作预算三个维度。只有身份、动作、接触、声音和时间关系共同满足，开篇的片段才算完成，一幅漂亮画面、一条自然声轨或一个汇总分都不能替代。

回到第一节的五项挑战：时间一致性由整段共同建模与自编码器压缩回答，同步由按时间对齐的条件接入和每步双向读取回答，几何一致性由共享三维表示、跨视角注意力和观测拟合回答，预算则通过压缩、级联和统一的验收口径变得可以比较。只有长时状态保持一项，本章的方法仅改善了片段之间的衔接。

仍未解决的部分同样清楚。本章的视频与音视频模型都只覆盖约五秒到三十秒的片段，长时身份与场景漂移需要持续保存状态；局部修改需要协调新旧内容的边界；联合音画已能学习共同事件，但口型和多声源的验证仍依赖人工；多视角创建的背面始终没有观测依据；流式回答要协调证据到达和语言播放。扩大主干或增加条件只提供更多表示和学习机会，究竟解决了哪个问题，仍要通过相应时间范围、输入条件和制作预算下的结果回答。

\paragraph{扩展阅读}
视频生成方向，Video Diffusion Models把图像扩散的U形网络分解为空间层和时间层，图像与视频联合训练，并用重建引导的条件采样延伸片段，可看作Lumiere一类时空网络的早期形式；Make-A-Video从文生图模型出发，只用无文字标注的视频学习运动，说明文字与外观可以先在图文数据上学好；Imagen Video用基础视频模型加多级空间与时间超分模型组成级联，正是Lumiere试图避免的时间级联结构\scite{vision-vdm2022,vision-makeavideo2022,vision-imagenvideo2022}Stable Video Diffusion在潜空间视频扩散上着重讨论数据筛选与图像预训练、视频预训练、高质量微调三阶段训练，可与Wan按阶段提高分辨率的训练安排对照；VideoPoet则把视频和音频离散成词元，用只含解码器的语言模型在同一序列中完成生成、编辑和配声等多种任务，代表扩散与流匹配之外的自回归路线\scite{vision-svd2023,vision-videopoet2024}

声音生成方向，AudioLDM用对比学习得到的音频—文字嵌入作为条件，在潜空间扩散生成声音，只接受文字条件，可用来理解MMAudio为何还要加入视觉同步特征；Diff-Foley先用音画对比预训练得到与时间对齐的视觉特征，再以它们为条件做潜扩散配声，与MMAudio同样把同步放进条件接入\scite{vision-audioldm2023,vision-difffoley2023}三维生成方向，Magic3D沿用分数蒸馏，先优化低分辨率神经辐射场，再转为带纹理的网格在更高分辨率上细化；ProlificDreamer把三维参数视为随机变量，提出变分分数蒸馏，针对DreamFusion式蒸馏常见的过饱和、过平滑和缺乏多样性\scite{vision-magic3d2023,vision-prolificdreamer2023}Zero-1-to-3以相对相机位姿为条件微调图像扩散模型，从单幅图像逐个合成新视角，各视角之间没有共同计算，可与Wonder3D在生成阶段协调多视角的做法对照；LRM则不逐物体优化，用Transformer从单幅图像直接前馈预测三平面表示，几秒内给出三维结果，省去分数蒸馏的逐例优化，前提是有大规模多视角数据用于训练\scite{vision-zero123-2023,vision-lrm2024}评测方面，EvalCrafter用一组提示词和多项客观指标评估视频生成模型，并按用户评分调整各指标的组合权重，可与VBench的维度划分相互参照\scite{vision-evalcrafter2024}
